\documentclass[12pt,a4paper]{ctexart}
\usepackage[utf8]{inputenc}
\usepackage{ctex}
\usepackage{amsmath,amssymb,amsfonts}
\usepackage{mathtools}
\usepackage{bm}
\usepackage{graphicx}
\usepackage{hyperref}
\usepackage{geometry}
\usepackage{fancyhdr}

\geometry{left=2.5cm,right=2.5cm,top=2.5cm,bottom=2.5cm}

\title{第三章\quad 机器人动力学}
\author{}
\date{}

\begin{document}

\maketitle

本章为机器人机械臂控制研究提供必要的背景知识。机械臂动力学方程采用二阶微分方程形式和多种状态变量形式推导。介绍动力学的一些重要性质。展示如何包含驱动机械臂的电机的动力学，这些电机可能是电动或液压马达。

\section{引言}

机器人学是一个涉及众多不同学科的复杂领域，如物理学、材料性质、静力学与动力学、电子学、控制理论、视觉、信号处理、计算机编程和制造技术。本书的主要兴趣是机器人机械臂的控制。本章的目的是研究控制机器人所需的动力学方程。

对于那些需要控制理论背景的读者，提供了第二章。对于那些需要机器人机械臂基础知识背景的读者，在附录A中考察了机器人机械臂的几何结构，涵盖了基本机械臂构型、运动学和逆运动学。在那里还回顾了机械臂雅可比矩阵，这对于在笛卡尔坐标系或工作空间坐标系中的控制至关重要，期望的机械臂轨迹通常首先在这些坐标系中指定。

机器人动力学在3.2节中推导。推导中使用了拉格朗日力学。在3.3节中，我们回顾了机械臂动力学方程的一些基本性质，这些性质对于后续章节中推导机器人控制方案至关重要。这些性质总结在表3.3.1中，在整本书中都会引用。

3.2节中的机械臂动力学采用二阶向量微分方程形式。在3.4节中，展示了将这种表述转换为状态变量描述的几种方法。状态变量描述是一阶向量微分方程，对于开发许多机械臂控制方案非常有用。本节使用了反馈线性化技术和哈密顿力学。

3.2节中的机器人机械臂动力学以关节空间坐标给出。在3.5节中，展示了一种非常通用的方法，用于以任何期望的坐标（包括笛卡尔坐标或工作空间坐标，以及相机坐标系或参考系的坐标）获得机械臂动力学描述。

在3.6节中，分析了执行移动机器人机械臂连杆所需工作的电动或液压驱动器。展示了如何将驱动器的动力学模型并入机械臂动力学中，以提供机械臂加驱动器系统的完整动力学描述。这最终使我们能够进入下一章讨论机器人机械臂控制设计。

\section{拉格朗日-欧拉动力学}

为了控制设计目的，需要一个揭示系统动力学行为的数学模型。因此，在本节中推导机器人机械臂的动力学方程。我们的方法是推导机械臂的动能和势能，然后使用拉格朗日运动方程。

在本节中，我们忽略驱动机器人机械臂的电动或液压电机的动力学；驱动器动力学在3.6节中讨论。

\subsection{力、惯量与能量}

让我们回顾物理学中的一些基本概念，这将使我们能够更好地理解机械臂动力学[Marion 1965]。在本小节中，我们使用粗体表示向量，普通字体表示其大小。

质量为$m$的物体以角速度$\omega$绕半径为$r$的点旋转时，其向心力为
\begin{equation}
F_{cen} = -m\omega \times (\omega \times r) = m\omega^2 r
\end{equation}

见图3.2.1。线速度为
\begin{equation}
v = \omega \times r
\end{equation}
在这种情况下意味着简单地$v = \omega r$。

想象一个球体（即地球）以其中心为轴以角速度$\omega_0$旋转。见图3.2.2。在球面上以速度$v$移动的质量为$m$的物体所受的科里奥利力为
\begin{equation}
F_{cor} = -2m\omega_0 \times v
\end{equation}

使用右手螺旋定则（即如果手指从$\omega_0$旋转到$v$，拇指指向$\omega_0 \times v$的方向），我们看到，在图中，科里奥利力使$m$向右偏转。

在低气压天气系统中，气团向低压中心移动。科里奥利力负责将气团向右偏转，从而引起称为气旋流的逆时针环流。结果是飓风的旋转运动。对图3.2.2的简要检查表明，在南半球，$F_{cor}$使移动的物体向左偏转，因此低压系统将具有顺时针风向运动。

由于$v = \omega \times r$和$\omega \times (\omega \times r) = \omega(\omega \cdot r) - r(\omega \cdot \omega)$，我们可以写出
\begin{equation}
F_{cen} = -m\omega \times (\omega \times r) = m\omega^2 r
\end{equation}

重要的是要注意，向心力涉及单个角速度的平方，而科里奥利力涉及两个不同角速度的乘积。

以线速度$v$移动的质量的动能为
\begin{equation}
K = \frac{1}{2}mv^2
\end{equation}

图3.2.1中质量的旋转动能为
\begin{equation}
K_{rot} = \frac{1}{2}I\omega^2
\end{equation}
其中转动惯量为
\begin{equation}
I = \int_V \rho(r)r^2 dr
\end{equation}
$\rho(r)$是体积中半径$r$处的质量分布。在所示的简单情况下，$m$是质点，这变为
\begin{equation}
I = mr^2
\end{equation}
因此，
\begin{equation}
K_{rot} = \frac{1}{2}mr^2\omega^2
\end{equation}

质量为$m$的物体在重力加速度为常数$g$的引力场中高度$h$处的势能为
\begin{equation}
P = mgh
\end{equation}
对应零势能的原点可以任意选择，因为只有势能差在物理力方面才有意义。

以速度$v$移动的质量$m$的动量为
\begin{equation}
p = mv
\end{equation}

质量$m$相对于原点的角动量为
\begin{equation}
P_{ang} = r \times p
\end{equation}

力$F$相对于同一原点的扭矩或力矩定义为
\begin{equation}
N = r \times F
\end{equation}

\subsection{拉格朗日运动方程}

对于保守系统，拉格朗日运动方程由下式给出[Marion 1965]
\begin{equation}
\frac{d}{dt}\frac{\partial L}{\partial \dot{q}} - \frac{\partial L}{\partial q} = \tau
\end{equation}
其中$q$是广义坐标$q_i$的$n$维向量，$\tau$是广义力$\tau_i$的$n$维向量，拉格朗日量是动能与势能之差
\begin{equation}
L = K - P
\end{equation}

在我们的使用中，$q$将是关节变量向量，由关节角$\theta_i$（以度或弧度表示）和关节偏移$d_i$（以米为单位）组成。那么$\tau$是一个向量，其分量为$n_i$，对应关节角的扭矩（牛顿-米），以及$f_i$，对应关节偏移的力（牛顿）。注意我们用大写字母表示$\tau$的分量。

我们将使用拉格朗日方程推导一般机器人机械臂动力学。让我们先通过考虑一些例子来感受一下这个过程。

\textbf{例3.2-1：二连杆极坐标机械臂的动力学}

二连杆平面旋转/平移（RP）机械臂的运动学在例A.2-3中给出。为了确定其动力学，请检查图3.2.3，其中关节变量和关节速度向量为
\begin{equation}
q = \begin{bmatrix} \theta \\ r \end{bmatrix}, \quad \dot{q} = \begin{bmatrix} \dot{\theta} \\ \dot{r} \end{bmatrix}
\end{equation}

相应的广义力向量为
\begin{equation}
\tau = \begin{bmatrix} n \\ f \end{bmatrix}
\end{equation}
其中$n$为扭矩，$f$为力。扭矩$n$和力$f$可由电机或液压驱动器提供。我们在3.6节讨论驱动器的动力学。

为了确定机械臂动力学，我们现在必须计算拉格朗日方程所需的量。

\textbf{a. 动能和势能}

由于角运动$\dot{\theta}$和线运动$\dot{r}$产生的总动能为
\begin{equation}
K = \frac{1}{2}m_1r^2\dot{\theta}^2 + \frac{1}{2}m_2(\dot{r}^2 + r^2\dot{\theta}^2)
\end{equation}
势能为
\begin{equation}
P = m_1gr\sin\theta + m_2gr\sin\theta
\end{equation}

\textbf{b. 拉格朗日方程}

拉格朗日量为
\begin{equation}
L = K - P = \frac{1}{2}m_1r^2\dot{\theta}^2 + \frac{1}{2}m_2(\dot{r}^2 + r^2\dot{\theta}^2) - m_1gr\sin\theta - m_2gr\sin\theta
\end{equation}

现在我们得到
\begin{align}
\frac{\partial L}{\partial \dot{q}} &= \begin{bmatrix} (m_1+m_2)r^2\dot{\theta} \\ m_2\dot{r} \end{bmatrix} \\
\frac{d}{dt}\frac{\partial L}{\partial \dot{q}} &= \begin{bmatrix} (m_1+m_2)r^2\ddot{\theta} + 2(m_1+m_2)r\dot{r}\dot{\theta} \\ m_2\ddot{r} \end{bmatrix} \\
\frac{\partial L}{\partial q} &= \begin{bmatrix} -(m_1+m_2)gr\cos\theta + (m_1+m_2)r\dot{\theta}^2 \\ (m_1+m_2)\dot{\theta}^2r - (m_1+m_2)g\sin\theta \end{bmatrix}
\end{align}

因此，(3.2.14)表明机械臂动力学方程为
\begin{align}
(m_1+m_2)r^2\ddot{\theta} + 2(m_1+m_2)r\dot{r}\dot{\theta} + (m_1+m_2)gr\cos\theta &= n \\
m_2\ddot{r} - (m_1+m_2)r\dot{\theta}^2 + (m_1+m_2)g\sin\theta &= f
\end{align}

这是一组耦合的非线性微分方程，描述了给定控制输入扭矩$n(t)$和力$f(t)$时的运动$q(t) = [\theta(t) \ r(t)]^T$。我们将在第四章展示如何通过计算机仿真确定给定控制输入$n(t)$和$f(t)$时的$q(t)$。

根据我们对力和惯量的讨论，很容易识别动力学方程中的各项。每个方程中的第一项是涉及质量和惯量的加速度项。(9)中的第二项是科里奥利项，而(10)中的第二项是向心项。第三项是重力项。

\textbf{c. 机械臂动力学}

通过使用向量，机械臂方程可以写成方便的形式。事实上，注意
\begin{equation}
\begin{bmatrix} (m_1+m_2)r^2 & 0 \\ 0 & m_2 \end{bmatrix} \begin{bmatrix} \ddot{\theta} \\ \ddot{r} \end{bmatrix} + \begin{bmatrix} 2(m_1+m_2)r\dot{r}\dot{\theta} \\ -(m_1+m_2)r\dot{\theta}^2 \end{bmatrix} + \begin{bmatrix} (m_1+m_2)gr\cos\theta \\ (m_1+m_2)g\sin\theta \end{bmatrix} = \begin{bmatrix} n \\ f \end{bmatrix}
\end{equation}

我们将此向量方程符号化为
\begin{equation}
M(q)\ddot{q} + V(q,\dot{q}) + G(q) = \tau
\end{equation}

注意，确实，惯量矩阵$M(q)$是$q$的函数（即$\theta$和$r$的函数），科里奥利/向心向量$V(q,\dot{q})$是$q$和$\dot{q}$的函数，重力向量$G(q)$是$q$的函数。

\textbf{例3.2-2：二连杆平面肘形机械臂的动力学}

在例A.2-2中给出了二连杆平面RR机械臂的运动学。为了确定其动力学，请检查图3.2.4，其中我们假设连杆质量集中在连杆末端。关节变量为
\begin{equation}
q = [\theta_1 \ \theta_2]^T
\end{equation}
广义力向量为
\begin{equation}
\tau = \begin{bmatrix} \tau_1 \\ \tau_2 \end{bmatrix}
\end{equation}
其中$\tau_1$和$\tau_2$是由驱动器提供的扭矩。

\textbf{a. 动能和势能}

对于连杆1，动能和势能为
\begin{align}
K_1 &= \frac{1}{2}m_1a_1^2\dot{\theta}_1^2 \\
P_1 &= m_1ga_1\sin\theta_1
\end{align}

对于连杆2，我们有
\begin{align}
x_2 &= a_1\cos\theta_1 + a_2\cos(\theta_1+\theta_2) \\
y_2 &= a_1\sin\theta_1 + a_2\sin(\theta_1+\theta_2) \\
\dot{x}_2 &= -a_1\sin\theta_1\dot{\theta}_1 - a_2\sin(\theta_1+\theta_2)(\dot{\theta}_1+\dot{\theta}_2) \\
\dot{y}_2 &= a_1\cos\theta_1\dot{\theta}_1 + a_2\cos(\theta_1+\theta_2)(\dot{\theta}_1+\dot{\theta}_2)
\end{align}

因此，速度平方为
\begin{equation}
v_2^2 = \dot{x}_2^2 + \dot{y}_2^2 = a_1^2\dot{\theta}_1^2 + a_2^2(\dot{\theta}_1+\dot{\theta}_2)^2 + 2a_1a_2(\dot{\theta}_1+\dot{\theta}_2)\dot{\theta}_1\cos\theta_2
\end{equation}

因此，连杆2的动能为
\begin{equation}
K_2 = \frac{1}{2}m_2v_2^2 = \frac{1}{2}m_2a_1^2\dot{\theta}_1^2 + \frac{1}{2}m_2a_2^2(\dot{\theta}_1+\dot{\theta}_2)^2 + m_2a_1a_2(\dot{\theta}_1+\dot{\theta}_2)\dot{\theta}_1\cos\theta_2
\end{equation}

连杆2的势能为
\begin{equation}
P_2 = m_2gy_2 = m_2g[a_1\sin\theta_1 + a_2\sin(\theta_1+\theta_2)]
\end{equation}

\textbf{b. 拉格朗日方程}

整个机械臂的拉格朗日量为
\begin{align}
L &= K - P = K_1 + K_2 - P_1 - P_2 \\
&= \frac{1}{2}(m_1+m_2)a_1^2\dot{\theta}_1^2 + \frac{1}{2}m_2a_2^2(\dot{\theta}_1+\dot{\theta}_2)^2 \\
&\quad + m_2a_1a_2(\dot{\theta}_1+\dot{\theta}_2)\dot{\theta}_1\cos\theta_2 - (m_1+m_2)ga_1\sin\theta_1 \\
&\quad - m_2ga_2\sin(\theta_1+\theta_2)
\end{align}

(3.2.14)所需的项为
\begin{align}
\frac{\partial L}{\partial \dot{q}} &= \begin{bmatrix} (m_1+m_2)a_1^2\dot{\theta}_1 + m_2a_2^2(\dot{\theta}_1+\dot{\theta}_2) + m_2a_1a_2(2\dot{\theta}_1+\dot{\theta}_2)\cos\theta_2 \\ m_2a_2^2(\dot{\theta}_1+\dot{\theta}_2) + m_2a_1a_2\dot{\theta}_1\cos\theta_2 \end{bmatrix} \\
\frac{d}{dt}\frac{\partial L}{\partial \dot{q}} &= \begin{bmatrix} (m_1+m_2)a_1^2\ddot{\theta}_1 + m_2a_2^2(\ddot{\theta}_1+\ddot{\theta}_2) + m_2a_1a_2(2\ddot{\theta}_1+\ddot{\theta}_2)\cos\theta_2 - m_2a_1a_2(2\dot{\theta}_1+\dot{\theta}_2)\dot{\theta}_2\sin\theta_2 \\ m_2a_2^2(\ddot{\theta}_1+\ddot{\theta}_2) + m_2a_1a_2\ddot{\theta}_1\cos\theta_2 - m_2a_1a_2\dot{\theta}_1\dot{\theta}_2\sin\theta_2 \end{bmatrix} \\
\frac{\partial L}{\partial q} &= \begin{bmatrix} -(m_1+m_2)ga_1\cos\theta_1 - m_2ga_2\cos(\theta_1+\theta_2) \\ -m_2a_1a_2(\dot{\theta}_1+\dot{\theta}_2)\dot{\theta}_1\sin\theta_2 - m_2ga_2\cos(\theta_1+\theta_2) \end{bmatrix}
\end{align}

最后，根据拉格朗日方程，机械臂动力学由两个耦合的非线性微分方程给出
\begin{align}
&[(m_1+m_2)a_1^2 + m_2a_2^2 + 2m_2a_1a_2\cos\theta_2]\ddot{\theta}_1 + [m_2a_2^2 + m_2a_1a_2\cos\theta_2]\ddot{\theta}_2 \\
&\quad - m_2a_1a_2(2\dot{\theta}_1+\dot{\theta}_2)\dot{\theta}_2\sin\theta_2 + (m_1+m_2)ga_1\cos\theta_1 + m_2ga_2\cos(\theta_1+\theta_2) = \tau_1 \nonumber \\
&[m_2a_2^2 + m_2a_1a_2\cos\theta_2]\ddot{\theta}_1 + m_2a_2^2\ddot{\theta}_2 + m_2a_1a_2\dot{\theta}_1^2\sin\theta_2 + m_2ga_2\cos(\theta_1+\theta_2) = \tau_2 \nonumber
\end{align}

\textbf{c. 机械臂动力学}

将机械臂动力学写成向量形式得到
\begin{equation}
\begin{bmatrix} (m_1+m_2)a_1^2 + m_2a_2^2 + 2m_2a_1a_2\cos\theta_2 & m_2a_2^2 + m_2a_1a_2\cos\theta_2 \\ m_2a_2^2 + m_2a_1a_2\cos\theta_2 & m_2a_2^2 \end{bmatrix} \begin{bmatrix} \ddot{\theta}_1 \\ \ddot{\theta}_2 \end{bmatrix}
\end{equation}
\begin{equation}
+ \begin{bmatrix} -m_2a_1a_2(2\dot{\theta}_1+\dot{\theta}_2)\dot{\theta}_2\sin\theta_2 \\ m_2a_1a_2\dot{\theta}_1^2\sin\theta_2 \end{bmatrix} + \begin{bmatrix} (m_1+m_2)ga_1\cos\theta_1 + m_2ga_2\cos(\theta_1+\theta_2) \\ m_2ga_2\cos(\theta_1+\theta_2) \end{bmatrix} = \begin{bmatrix} \tau_1 \\ \tau_2 \end{bmatrix}
\end{equation}
其中
\begin{equation}
M(q) = \begin{bmatrix} (m_1+m_2)a_1^2 + m_2a_2^2 + 2m_2a_1a_2\cos\theta_2 & m_2a_2^2 + m_2a_1a_2\cos\theta_2 \\ m_2a_2^2 + m_2a_1a_2\cos\theta_2 & m_2a_2^2 \end{bmatrix}
\end{equation}

这些机械臂动力学采用标准形式
\begin{equation}
M(q)\ddot{q} + V(q,\dot{q}) + G(q) = \tau
\end{equation}
其中$M(q)$为惯量矩阵，$V(q,\dot{q})$为科里奥利/向心向量，$G(q)$为重力向量。注意$M(q)$是对称的。

\textbf{习题3.2-3：三连杆圆柱形机械臂的动力学}

我们在例A.2-1中研究三连杆圆柱形机械臂的运动学。在图3.2.5中，关节变量向量为
\begin{equation}
q = [\theta \ h \ r]^T
\end{equation}

证明机械臂动力学由下式给出
\begin{equation}
\begin{bmatrix} J+m_1r^2+m_2r^2 & 0 & 0 \\ 0 & m_1+m_2 & 0 \\ 0 & 0 & m_2 \end{bmatrix} \begin{bmatrix} \ddot{\theta} \\ \ddot{h} \\ \ddot{r} \end{bmatrix} + \begin{bmatrix} 2(m_1+m_2)r\dot{r}\dot{\theta} \\ 0 \\ -(m_1+m_2)r\dot{\theta}^2 \end{bmatrix} + \begin{bmatrix} 0 \\ (m_1+m_2)g \\ 0 \end{bmatrix} = \begin{bmatrix} n \\ f_h \\ f_r \end{bmatrix}
\end{equation}
其中$J$是基座连杆的惯量，力向量为
\begin{equation}
\tau = \begin{bmatrix} n \\ f_h \\
 f_r \end{bmatrix}
\end{equation}

\subsection{推导操作器动力学}

我们在几个例子中展示了如何应用拉格朗日方程计算任何给定机器人机械臂的动力学方程。在例子中，我们发现的动力学总是具有特殊形式
\begin{equation}
M(q)\ddot{q} + V(q,\dot{q}) + G(q) = \tau
\end{equation}
其中$q$为关节变量向量，$\tau$为广义力/扭矩向量。在本小节中，我们推导一般机器人机械臂的动力学。它们将采用这种相同的形式。

为了获得一般机器人机械臂动力学方程，我们确定机械臂动能和势能，然后拉格朗日量，最后代入拉格朗日方程(3.2.14)获得最终结果[Paul 1981, Lee et al. 1983, Asada and Slotine 1986, Spong and Vidyasagar 1989]。

\textbf{机械臂动能}

给定连杆$i$上以相对于连接到该连杆的坐标系$i$的坐标$^ir$表示的一个点，该点在基坐标系中的坐标为
\begin{equation}
\begin{bmatrix} x \\ y \\ z \\ 1 \end{bmatrix} = T_i \, {^ir}
\end{equation}
其中$T_i$是附录A中定义的$4 \times 4$齐次变换矩阵。注意$T_i$是关节变量$q_1, q_2, \ldots, q_i$的函数。因此，该点在基坐标系中的速度为
\begin{equation}
\begin{bmatrix} \dot{x} \\ \dot{y} \\ \dot{z} \\ 0 \end{bmatrix} = \sum_{j=1}^{i} \frac{\partial T_i}{\partial q_j}\dot{q}_j \, {^ir}
\end{equation}

由于当$j > i$时，$\partial T_i/\partial q_j = 0$，我们可以将上限求和替换为$n$，即连杆数。如果已知机械臂矩阵$T_i$，则可以计算$4 \times 4$矩阵$\partial T_i/\partial q_j$。

以速度$v = [v_x \ v_y \ v_z]^T$运动的无穷小质量$dm$的动能为
\begin{equation}
dK_i = \frac{1}{2}(\dot{x}^2 + \dot{y}^2 + \dot{z}^2)dm = \frac{1}{2}\text{trace}(\dot{T}_i \, {^ir} \, {^ir}^T \dot{T}_i^T)dm
\end{equation}

因此，连杆$i$的总动能为
\begin{equation}
K_i = \int_{\text{link }i} dK_i = \frac{1}{2}\text{trace}\left[\sum_{j=1}^{i}\sum_{k=1}^{i}\frac{\partial T_i}{\partial q_j}\dot{q}_j\left(\int_{\text{link }i}{^ir} \, {^ir}^T dm\right)\frac{\partial T_i^T}{\partial q_k}\dot{q}_k\right]
\end{equation}

将$dK_i$的表达式代入，我们可以将积分移到求和内。然后，为连杆$i$定义$4 \times 4$伪惯量矩阵为
\begin{equation}
M_i = \int_{\text{link }i}{^ir} \, {^ir}^T dm
\end{equation}

我们可以将连杆$i$的动能写为
\begin{equation}
K_i = \frac{1}{2}\sum_{j=1}^{i}\sum_{k=1}^{i}\text{trace}\left(\frac{\partial T_i}{\partial q_j}M_i\frac{\partial T_i^T}{\partial q_k}\right)\dot{q}_j\dot{q}_k = \frac{1}{2}q^T M(q) \dot{q}
\end{equation}

让我们继续寻找机械臂总动能之前简要讨论伪惯量矩阵。令
\begin{equation}
{^ir} = [x \ y \ z \ 1]^T
\end{equation}
为无穷小质量$dm$在坐标系$i$中的坐标。那么，展开(3.2.20)得到
\begin{equation}
M_i = \begin{bmatrix} \int x^2 dm & \int xy dm & \int xz dm & \int x dm \\ \int xy dm & \int y^2 dm & \int yz dm & \int y dm \\ \int xz dm & \int yz dm & \int z^2 dm & \int z dm \\ \int x dm & \int y dm & \int z dm & \int dm \end{bmatrix}
\end{equation}
其中积分在连杆$i$的体积上进行。这是一个对每个连杆计算一次的常数矩阵。它取决于连杆$i$的几何形状和质量分布。实际上，根据连杆$i$的惯性矩
\begin{align}
I_{xx} &= \int(y^2+z^2)dm \\ I_{yy} &= \int(x^2+z^2)dm \\ I_{zz} &= \int(x^2+y^2)dm
\end{align}
惯性积
\begin{align}
I_{xy} &= \int xy dm \\ I_{xz} &= \int xz dm \\ I_{yz} &= \int yz dm
\end{align}
和一阶矩
\begin{align}
mx &= \int x dm = mx_c \\ my &= \int y dm = my_c \\ mz &= \int z dm = mz_c
\end{align}
其中$m$为连杆$i$的总质量，
\begin{equation}
x_c = \frac{1}{m}\int x dm
\end{equation}
是连杆$i$质心在坐标系$i$中的坐标，我们可以写出
\begin{equation}
M_i = \begin{bmatrix} \frac{1}{2}(-I_{xx}+I_{yy}+I_{zz}) & I_{xy} & I_{xz} & mx_c \\ I_{xy} & \frac{1}{2}(I_{xx}-I_{yy}+I_{zz}) & I_{yz} & my_c \\ I_{xz} & I_{yz} & \frac{1}{2}(I_{xx}+I_{yy}-I_{zz}) & mz_c \\ mx_c & my_c & mz_c & m \end{bmatrix}
\end{equation}

这些量要么在机械臂制造商的规格表中列出，要么可以从那里列出的量计算得出。

回到我们的推导，机械臂总动能可以写为
\begin{equation}
K = \sum_{i=1}^{n}K_i = \frac{1}{2}\sum_{i=1}^{n}\sum_{j=1}^{i}\sum_{k=1}^{i}\text{trace}\left(\frac{\partial T_i}{\partial q_j}M_i\frac{\partial T_i^T}{\partial q_k}\right)\dot{q}_j\dot{q}_k
\end{equation}

由于矩阵和的迹是各个迹的和，我们可以交换求和和迹算子得到
\begin{equation}
K = \frac{1}{2}\sum_{j=1}^{n}\sum_{k=1}^{n}m_{jk}\dot{q}_j\dot{q}_k = \frac{1}{2}\dot{q}^T M(q) \dot{q}
\end{equation}
或
\begin{equation}
K = \frac{1}{2}\dot{q}^T M(q) \dot{q}
\end{equation}
其中$n \times n$机械臂惯量矩阵$M(q)$的元素定义为
\begin{equation}
m_{jk} = \sum_{i=\max(j,k)}^{n}\text{trace}\left(\frac{\partial T_i}{\partial q_j}M_i\frac{\partial T_i^T}{\partial q_k}\right)
\end{equation}

由于当$j > i$时，$\partial T_i/\partial q_j = 0$，我们可以更有效地将其写为
\begin{equation}
m_{jk} = \sum_{i=\max(j,k)}^{n}\text{trace}\left(\frac{\partial T_i}{\partial q_j}M_i\frac{\partial T_i^T}{\partial q_k}\right)
\end{equation}

方程(3.2.29)正是我们一直在寻求的；它提供了以已知量和关节变量$q$表示的机械臂动能的方便表达式。由于$m_{jk} = m_{kj}$，惯量矩阵$M(q)$是对称的。由于动能是正的，只有当广义速度为零时才消失，惯量矩阵$M(q)$也是正定的。注意动能取决于$q$和$\dot{q}$。

\textbf{机械臂势能}

如果连杆$i$具有质量$m_i$和以其坐标系$i$的坐标表示的质心，该连杆的势能为
\begin{equation}
P_i = -m_i g^T T_i {^ir_i}
\end{equation}
其中重力向量以基坐标表示为
\begin{equation}
g = \begin{bmatrix} g_x \\ g_y \\ g_z \\ 0 \end{bmatrix}
\end{equation}

如果机械臂是水平的，在海平面上，基座$z$轴垂直向上，那么
\begin{equation}
g = \begin{bmatrix} 0 \\ 0 \\ -9.8 \\ 0 \end{bmatrix} \text{m/s}^2
\end{equation}

因此，机械臂总势能为
\begin{equation}
P = \sum_{i=1}^{n}P_i = -\sum_{i=1}^{n}m_i g^T T_i {^ir_i}
\end{equation}
注意$P$仅取决于关节变量$q$，而不取决于关节速度$\dot{q}$。

注意$-m_i g^T T_i$是连杆$i$伪惯量矩阵$M_i$的最后一列，我们可以写出
\begin{equation}
P = -\sum_{i=1}^{n}m_i g^T T_i {^ir_i}
\end{equation}
其中$e_4$是$4 \times 4$单位矩阵的最后一列（即$e_4 = [0 \ 0 \ 0 \ 1]^T$）。

\textbf{拉格朗日方程}

机械臂拉格朗日量为
\begin{equation}
L = K - P = \frac{1}{2}\dot{q}^T M(q)\dot{q} - P(q)
\end{equation}

动能是关节速度向量的二次函数，势能独立于$\dot{q}$是一个基本性质。

拉格朗日方程(3.2.14)现在所需的项由下式给出
\begin{align}
\frac{\partial L}{\partial \dot{q}} &= M(q)\dot{q} \\
\frac{d}{dt}\frac{\partial L}{\partial \dot{q}} &= M(q)\ddot{q} + \dot{M}(q)\dot{q} \\
\frac{\partial L}{\partial q} &= \frac{1}{2}\frac{\partial}{\partial q}(\dot{q}^T M(q)\dot{q}) - \frac{\partial P(q)}{\partial q}
\end{align}

因此，机械臂动力学方程为
\begin{equation}
M(q)\ddot{q} + \dot{M}(q)\dot{q} - \frac{1}{2}\frac{\partial}{\partial q}(\dot{q}^T M(q)\dot{q}) + \frac{\partial P(q)}{\partial q} = \tau
\end{equation}

定义科里奥利/向心向量
\begin{equation}
V(q,\dot{q}) = \dot{M}(q)\dot{q} - \frac{1}{2}\frac{\partial}{\partial q}(\dot{q}^T M(q)\dot{q})
\end{equation}
和重力向量
\begin{equation}
G(q) = \frac{\partial P(q)}{\partial q}
\end{equation}

我们可以写出
\begin{equation}
M(q)\ddot{q} + V(q,\dot{q}) + G(q) = \tau
\end{equation}
这正是我们一直在寻求的机器人动力学方程的最终形式。

$M(q)$元素相对于旋转关节变量$q_i = \theta_i$的单位是kg-$\text{m}^2$。$M(q)$元素相对于平移关节变量$q_i = d_i$的单位是千克。$V(q,\dot{q})$和$G(q)$元素相对于旋转关节变量的单位是kg-$\text{m}^2/\text{s}^2$。$V(q,\dot{q})$和$G(q)$元素相对于平移关节变量的单位是kg-m/$\text{s}^2$。

\section{机器人方程的结构与性质}

在本节中，我们研究动力学机械臂方程的详细结构和性质，因为这种结构应该反映在控制律的形式中。如果在设计阶段结合机械臂的已知性质，控制器会更简单更有效。

实际上，机器人机械臂总是受摩擦和干扰的影响。因此，我们将我们刚刚推导出的机械臂模型推广，将机械臂动力学写为
\begin{equation}
M(q)\ddot{q} + V(q,\dot{q}) + F_v\dot{q} + F_d(\dot{q}) + G(q) + \tau_d = \tau
\end{equation}
其中$q$为关节变量$n$维向量，$\tau$为广义力的$n$维向量。$M(q)$是惯量矩阵，$V(q,\dot{q})$是科里奥利/向心向量，$G(q)$是重力向量。我们添加了一个摩擦项
\begin{equation}
F(\dot{q}) = F_v\dot{q} + F_d(\dot{q})
\end{equation}
其中$F_v$是粘性摩擦系数矩阵，$F_d$是动态摩擦项。还添加了一个干扰$\tau_d$，它可以表示，例如，任何建模不准确的动力学。

摩擦不是一个容易建模的项，实际上，可能是机械臂动力学模型中最难描述的项。关于摩擦的更多讨论可以在[Schilling 1990]中找到。

我们有时将机械臂动力学写为
\begin{equation}
M(q)\ddot{q} + N(q,\dot{q}) = \tau
\end{equation}
其中
\begin{equation}
N(q,\dot{q}) = V(q,\dot{q}) + F(\dot{q}) + G(q) + \tau_d
\end{equation}
表示非线性项。

让我们检查机器人动力学方程中每一项的结构和性质。这项研究将为我们提供很大的洞察力，我们在后续章节中推导机器人控制方案时会用到。我们在表3.3.1中总结了发现的性质。当我们发展每个性质时，参考例3.2.1到3.2.3以验证这些性质确实在那里成立是值得的。在本节末尾的例3.3.1中，我们说明了二连杆平面肘形机械臂的几个性质。

\begin{table}[htbp]
\centering
\caption{机器人方程及其性质}
\begin{tabular}{|l|l|}
\hline
\textbf{机器人方程} & \\
\hline
动力学 & $M(q)\ddot{q} + V(q,\dot{q}) + F_v\dot{q} + F_d(\dot{q}) + G(q) + \tau_d = \tau$ \\
\hline
状态空间 & $\dot{x} = Ax + Bu$，其中$u = \tau$，$x = [q^T \ \dot{q}^T]^T$ \\
\hline
\textbf{性质} & \\
\hline
惯量矩阵 & 对称、正定的有界矩阵 \\
\hline
科里奥利/向心 & 包含$\dot{q}$的二次项 \\
\hline
摩擦 & $F(\dot{q}) = F_v\dot{q} + F_d(\dot{q})$ \\
\hline
干扰 & $||\tau_d|| \leq d$ \\
\hline
线性参数化 & $M(q)\ddot{q} + V(q,\dot{q}) + F_v\dot{q} + F_d(\dot{q}) + G(q) = W(q,\dot{q},\ddot{q})\phi$ \\
\hline
斜对称 & $S(q,\dot{q}) = \dot{M}(q) - 2V_m(q,\dot{q})$是斜对称的 \\
\hline
无源性 & $\int_0^T \dot{q}^T\tau dt \geq -\gamma^2$ \\
\hline
能量守恒 & $\dot{K} = \dot{q}^T(\tau - G)$ \\
\hline
\end{tabular}
\end{table}

\subsection{惯量矩阵性质}

正如我们所见，$M(q)$是对称且正定的。实际上，机械臂动能为
\begin{equation}
K = \frac{1}{2}\dot{q}^T M(q) \dot{q}
\end{equation}

下一小节给出了$M$的一些表达式。

$M(q)$的另一个重要性质是它上下有界。即，
\begin{equation}
\mu_1 I \leq M(q) \leq \mu_2 I
\end{equation}
其中$\mu_1$和$\mu_2$是可以为任何给定机械臂计算的标量（见例3.3.1）。当我们说$\mu_1 I \leq M(q)$，例如，我们的意思是$(M(q) - \mu_1 I)$是半正定的。即，
\begin{equation}
x^T(M - \mu_1 I)x \geq 0
\end{equation}
对所有$x \in R^n$成立。

同样，惯量矩阵的逆是有界的，因为
\begin{equation}
\frac{1}{\mu_2} I \leq M^{-1}(q) \leq \frac{1}{\mu_1} I
\end{equation}

如果机械臂是旋转关节的，界限$\mu_1$和$\mu_2$是常数，因为$q$只通过$\sin$和$\cos$项出现在$M(q)$中，它们的大小被1限制（见例3.2.2和3.3.1）。另一方面，如果机械臂有平移关节，那么$\mu_1$和$\mu_2$可能是$q$的标量函数。见例3.2.1，其中$M(q)$被$\mu_2 = mr^2$（如果$r > 1$）上界限制。

惯量矩阵的有界性也可以表示为
\begin{equation}
m_1 \leq ||M(q)|| \leq m_2
\end{equation}
其中任何诱导矩阵范数可以用来定义正标量$m_1$和$m_2$。

\subsection{科里奥利/向心项性质}

一瞥(3.2.42)就会发现一个问题，如果不理解，会使机器人动力学的研究变得混乱。这个$V(q,\dot{q})$项的简化将需要对矩阵[即$M(q)$]关于$n$维向量$q$求导。然而，这样的导数不是矩阵，而是三阶张量——即，它们必须用三个指标表示，而不是两个。有几种方法可以解决这个问题，涉及几个新量的定义。

\textbf{$V(q,\dot{q})$的克罗内克积分析}

让我们首先从克罗内克积[Brewer 1978]的角度检查项$V(q,\dot{q})$，克罗内克积对两个矩阵$A$、$B$定义为
\begin{equation}
A \otimes B = [a_{ij}B]
\end{equation}
其中$A$有元素$a_{ij}$，$[a_{ij}B]$表示由$p \times q$块$a_{ij}B$组成的$np \times mq$块矩阵。因此，对于$A$、$B$、$C$、$D$，我们有
\begin{equation}
(A \otimes B)(C \otimes D) = AC \otimes BD
\end{equation}

对于矩阵$A(q)$、$B(q)$，定义矩阵导数为
\begin{equation}
\frac{\partial A}{\partial q} = \begin{bmatrix} \frac{\partial A}{\partial q_1} & \cdots & \frac{\partial A}{\partial q_n} \end{bmatrix}
\end{equation}

然后我们可以证明乘积法则
\begin{equation}
\frac{\partial}{\partial q}(AB) = \frac{\partial A}{\partial q}(I_n \otimes B) + A\frac{\partial B}{\partial q}
\end{equation}
其中$I_n$是$n \times n$单位矩阵。

现在我们可以检查科里奥利/向心向量$V(q,\dot{q})$[Koditschek 1984, Gu and Loh 1988]。在(3.2.42)上使用(3.3.11)两次，我们可以获得
\begin{equation}
V(q,\dot{q}) = \frac{dM}{dt}\dot{q} - \frac{1}{2}(I_n \otimes \dot{q}^T)\frac{\partial M}{\partial q}\dot{q}
\end{equation}
或
\begin{equation}
V(q,\dot{q}) = [\dot{M}(q) - \frac{1}{2}U(q)]\dot{q}
\end{equation}
其中
\begin{align}
U(q) &= \begin{bmatrix} \frac{\partial M}{\partial q_1} & \cdots & \frac{\partial M}{\partial q_n} \end{bmatrix} \\
\dot{M} &= \sum_{j=1}^{n}\frac{\partial M}{\partial q_j}\dot{q}_j
\end{align}
矩阵系数由下式给出
\begin{equation}
\frac{\partial M}{\partial q_j} = \begin{bmatrix} \frac{\partial m_{1j}}{\partial q_1} & \cdots & \frac{\partial m_{1j}}{\partial q_n} \\ \vdots & \ddots & \vdots \\ \frac{\partial m_{nj}}{\partial q_1} & \cdots & \frac{\partial m_{nj}}{\partial q_n} \end{bmatrix}
\end{equation}

为了找到$V_{m1}$的等效表达式，注意
\begin{equation}
\frac{\partial}{\partial q}(\dot{q}^T M \dot{q}) = \frac{\partial}{\partial q}(\dot{q}^T M) \dot{q} + (I_n \otimes \dot{q}^T)\frac{\partial M^T}{\partial q}\dot{q}
\end{equation}

这可以写为
\begin{equation}
\frac{\partial}{\partial q}(\dot{q}^T M \dot{q}) = \left[\frac{\partial}{\partial q}(\dot{q}^T M) + (I_n \otimes \dot{q}^T)\frac{\partial M}{\partial q}\right]\dot{q}
\end{equation}
或
\begin{equation}
\frac{\partial}{\partial q}(\dot{q}^T M \dot{q}) = \left[\dot{q}^T\frac{\partial M}{\partial q} + (I_n \otimes \dot{q}^T)\frac{\partial M}{\partial q}\right]\dot{q}
\end{equation}
（注意$\partial M/\partial q_i$是对称的。）因此，
\begin{equation}
V_{m1}(q,\dot{q}) = \frac{1}{2}\left[\dot{M}(q) + (I_n \otimes \dot{q}^T)\frac{\partial M}{\partial q}\right]\dot{q}
\end{equation}

当使用适当的定义时，我们可以写出
\begin{equation}
V_{m1}(q,\dot{q}) = \frac{1}{2}[\dot{M}(q) + U^T(q)]\dot{q}
\end{equation}

也可以写出（见习题）
\begin{equation}
V_{m2}(q,\dot{q}) = \frac{1}{2}[\dot{M}(q) + U(q)]\dot{q}
\end{equation}

由于$V_v(\dot{q})$在$\dot{q}$中是线性的，因此$V(q,\dot{q})$在$\dot{q}$中是二次的。事实上，可以证明（见习题）
\begin{equation}
V(q,\dot{q}) = \begin{bmatrix} \dot{q}^T V_1(q)\dot{q} \\ \vdots \\ \dot{q}^T V_n(q)\dot{q} \end{bmatrix}
\end{equation}
对于$V_i(q)$的适当定义[Craig 1988]。确实，$V_i(q)$是对称的$n \times n$矩阵。

由于$V(q,\dot{q})$在$\dot{q}$中是二次的，它可以被$\dot{q}$的二次函数上界限制。即，
\begin{equation}
||V(q,\dot{q})|| \leq v_b(q)||\dot{q}||^2
\end{equation}
其中$v_b(q)$是已知的标量函数，$||\cdot||$是任何适当的范数。对于旋转机械臂，$v_b$是与$q$无关的常数。见例3.2.2和3.3.1，其中$\dot{q}$中的二次项乘以$\sin\theta_2$，其大小被1限制。另一方面，对于具有平移关节的机械臂，$v_b(q)$可能是$q$的函数；见例3.2.1和3.2.3，其中$V(q,\dot{q})$有一个在$\dot{q}$中的二次项乘以$r$。

为了帮助确定给定机器人机械臂的$v_b(q)$，注意$V(q,\dot{q}) = V_v(\dot{q})\dot{q}$，因此
\begin{equation}
||V(q,\dot{q})|| \leq ||V_v(q)|| \cdot ||\dot{q}||
\end{equation}
其中$||V_v(q)||$在(3.3.23)中定义。因此，对于旋转机械臂
\begin{equation}
v_b = \sup_{q} ||V_v(q)||
\end{equation}

我们可以注意到$\dot{q}^2 \equiv \dot{q} \otimes \dot{q}$，因此
\begin{equation}
\dot{q}^2 = \begin{bmatrix} \dot{q}_1^2 & \dot{q}_1\dot{q}_2 & \cdots & \dot{q}_1\dot{q}_n & \dot{q}_2\dot{q}_1 & \cdots & \dot{q}_n^2 \end{bmatrix}^T
\end{equation}
这是一个由$\dot{q}$分量所有可能乘积组成的$n^2$维向量。这允许我们证明
\begin{equation}
||\dot{q}||^2 = ||\dot{q}^2||
\end{equation}
在这个证明中，我们还需要恒等式
\begin{equation}
(A \otimes B)^T = A^T \otimes B^T
\end{equation}
对任何矩阵$A$和$B$成立。注意：$(A \otimes B)^T = A^T \otimes B^T$不成立。

现在我们可以使用这些各种恒等式来证明
\begin{equation}
V_m(q,\dot{q}) = V_{m1}(q) + V_{v1}(q)(I_n \otimes \dot{q}) = V_{m2}(q) - V_{v2}(q)(I_n \otimes \dot{q})
\end{equation}

因此
\begin{equation}
V(q,\dot{q}) = V_m(q)\dot{q} = V_{m1}(q,\dot{q})\dot{q} = V_{m2}(q,\dot{q})\dot{q}
\end{equation}
其中
\begin{align}
V_{m1}(q) &= \frac{1}{2}[\dot{M}(q) + U^T(q)] \\
V_{m2}(q) &= \frac{1}{2}[\dot{M}(q) + U(q)]
\end{align}
注意，一般来说，$V_{m1} \neq V_{m2}$。

就$M$和$U$而言，机械臂动力学可以写为
\begin{equation}
M(q)\ddot{q} + V_m(q)\dot{q} + G(q) = \tau
\end{equation}

在这一点上，我们可以证明一个恒等式，它在构建高级控制方案时非常有用。我们称之为斜对称性质；它表明$M(q)$的导数和科里奥利向量以非常特殊的方式相关。事实上，
\begin{equation}
\dot{M} - 2V_m = -(\dot{M} - 2V_m)^T
\end{equation}

由于矩阵减去其转置总是斜对称的，这是一个重要的恒等式。如果使用$V_{m1}$代替$V_{m2}$，这也成立。

重要的是要注意，(3.3.33)中的第一个等式成立是因为$\dot{q}$乘以$\dot{M} - 2V_m$。即，$\dot{q}^T(\dot{M} - 2V_m)\dot{q} = 0$，因此$(\dot{M} - 2V_m)$本身不一定是斜对称的。然而，可以定义一个矩阵$V_m$，使得
\begin{equation}
V(q,\dot{q}) = V_m(q,\dot{q})\dot{q}
\end{equation}
和
\begin{equation}
S(q,\dot{q}) \equiv \dot{M}(q) - 2V_m(q,\dot{q})
\end{equation}
是斜对称的，因此对所有$x \in R^n$，$x^T S x = 0$。确实，根据(3.3.13)和(3.3.27)，我们可以定义
\begin{equation}
V_m(q,\dot{q}) = \frac{1}{2}[\dot{M}(q) - U^T(q)]
\end{equation}
因为那么斜对称矩阵只不过是
\begin{equation}
S = \dot{M} - 2V_m = U - U^T
\end{equation}

这个$V_m$是在几个现代自适应和鲁棒控制算法中使用的标准定义，也是我们在本书其余部分要使用的定义。因此，我们将机械臂方程写为
\begin{equation}
M(q)\ddot{q} + V_m(q,\dot{q})\dot{q} + G(q) = \tau
\end{equation}
或
\begin{equation}
M(q)\ddot{q} + V(q,\dot{q}) + G(q) = \tau
\end{equation}

注意，可以将$V(q,\dot{q})$分成其科里奥利和向心分量为
\begin{equation}
V(q,\dot{q}) = V_{cor}(q) + V_{cen}(q,\dot{q})
\end{equation}
其中
\begin{equation}
V_{cor}(q) = \dot{M}(q)\dot{q} - V_{cen}(q,\dot{q})
\end{equation}
和$V_{cen}$是(3.3.26)中所有平方项$\dot{q}_i^2$被移除的$\dot{q}^2$[Craig 1988]（见习题）。

\textbf{$V(q,\dot{q})$的分量分析}

科里奥利/向心向量的克罗内克积分析的替代方法是基于$V(q,\dot{q})$标量分量的分析，这提供了额外的洞察力。

就惯量矩阵$M(q)$的分量$m_{kj}(q)$而言，我们可以将(3.2.38)-(3.2.40)分量写为
\begin{align}
\frac{\partial K}{\partial \dot{q}_k} &= \sum_{j}m_{kj}\dot{q}_j \\
\frac{d}{dt}\frac{\partial K}{\partial \dot{q}_k} &= \sum_{j}m_{kj}\ddot{q}_j + \sum_{i,j}\frac{\partial m_{kj}}{\partial q_i}\dot{q}_i\dot{q}_j \\
\frac{\partial K}{\partial q_k} &= \frac{1}{2}\sum_{i,j}\frac{\partial m_{ij}}{\partial q_k}\dot{q}_i\dot{q}_j
\end{align}
其中所有求和都是在关节数$n$上进行的。现在，拉格朗日方程表明机械臂动力学分量表示为
\begin{equation}
\sum_{j}m_{kj}\ddot{q}_j + \sum_{i,j}\left[\frac{\partial m_{kj}}{\partial q_i} - \frac{1}{2}\frac{\partial m_{ij}}{\partial q_k}\right]\dot{q}_i\dot{q}_j + \frac{\partial P}{\partial q_k} = \tau_k
\end{equation}
其中$n$是关节数。

通过交换求和顺序并利用对称性，
\begin{equation}
\sum_{i,j}\frac{\partial m_{kj}}{\partial q_i}\dot{q}_i\dot{q}_j = \frac{1}{2}\sum_{i,j}\left[\frac{\partial m_{kj}}{\partial q_i} + \frac{\partial m_{ki}}{\partial q_j}\right]\dot{q}_i\dot{q}_j
\end{equation}

因此，我们可以定义
\begin{equation}
v_{ijk} = \frac{1}{2}\left[\frac{\partial m_{kj}}{\partial q_i} + \frac{\partial m_{ki}}{\partial q_j} - \frac{\partial m_{ij}}{\partial q_k}\right]
\end{equation}
并将机械臂动力学写为
\begin{equation}
\sum_{j}m_{kj}\ddot{q}_j + \sum_{i,j}v_{ijk}\dot{q}_i\dot{q}_j + \frac{\partial P}{\partial q_k} = \tau_k
\end{equation}

$v_{ijk}$的循环对称性使我们能够推导出这个方程中对应于科里奥利/向心向量$V(q,\dot{q})$的第二项的重要性质。量$v_{ijk}$被称为（第一类）克里斯托费尔符号[Borisenko and Tarapov 1968]。

(3.3.36)中定义的矩阵$V_m$的分量$v_{kj}$由下式给出
\begin{equation}
v_{kj} = \sum_{i}v_{ijk}\dot{q}_i
\end{equation}

\subsection{重力、摩擦与干扰性质}

\textbf{重力项$G(q)$的性质}

根据(3.2.43)和(3.2.36)，
\begin{equation}
G(q) = \frac{\partial P}{\partial q} = -\sum_{i=1}^{n}m_i g^T \frac{\partial T_i}{\partial q}{^ir_i}
\end{equation}
因此使用(3.3.11)两次揭示
\begin{equation}
G(q) = -\sum_{i=1}^{n}m_i (I_n \otimes g^T)\frac{\partial T_i}{\partial q}{^ir_i}
\end{equation}

可以为任何给定机器人机械臂导出重力项的界限。因此
\begin{equation}
||G(q)|| \leq g_b(q)
\end{equation}
其中$||\cdot||$是任何适当的向量范数，$g_b$是可以为任何给定机械臂确定的标量函数（见例3.3.1）。对于旋转机械臂，$g_b$是与关节向量$q$无关的常数，但对于具有平移连杆的机械臂，$g_b$可能取决于$q$。见3.2节和例3.3.1中的例子来验证这些说法。

\textbf{摩擦项$F(\dot{q})$的性质}

机械臂方程(3.3.1)中的摩擦形式为
\begin{equation}
F(\dot{q}) = F_v\dot{q} + F_d(\dot{q})
\end{equation}
其中$F_v$是粘性摩擦系数矩阵，$F_d$是动态摩擦项。摩擦系数是为给定机械臂最难确定的参数之一，实际上，(3.3.52)仅代表其影响的近似数学模型。更多讨论见[Craig 1988, Schilling 1990]。

由于摩擦是局部效应，我们可以假设$F(\dot{q})$在关节之间是解耦的，因此
\begin{equation}
F(\dot{q}) = \begin{bmatrix} f_1(\dot{q}_1) \\ \vdots \\ f_n(\dot{q}_n) \end{bmatrix}
\end{equation}
其中$f_i(\cdot)$是可以为任何给定机械臂确定的已知标量函数。我们为将来使用定义了$\text{vec}\{\cdot\}$函数。

粘性摩擦通常可以假设具有形式
\begin{equation}
F_v\dot{q} = \text{diag}\{v_i\}\dot{q}
\end{equation}
其中$v_i$是已知常数系数。那么$F_v = \text{diag}\{v_i\}$，一个对角元素为$v_i$的对角矩阵。动态摩擦通常可以假设具有形式
\begin{equation}
F_d(\dot{q}) = \text{diag}\{k_i\}\text{sgn}(\dot{q})
\end{equation}
其中$k_i$是已知常数系数，符号函数对标量$x$定义为
\begin{equation}
\text{sgn}(x) = \begin{cases} +1, & x > 0 \\ -1, & x < 0 \\ 0, & x = 0 \end{cases}
\end{equation}

那么
\begin{equation}
F_d(\dot{q}) = K_d \text{sgn}(\dot{q})
\end{equation}
其中$K_d = \text{diag}\{k_i\}$是动态摩擦系数矩阵，向量$x$的符号函数定义为
\begin{equation}
\text{sgn}(x) = \text{vec}\{\text{sgn}(x_i)\} = \begin{bmatrix} \text{sgn}(x_1) \\ \vdots \\ \text{sgn}(x_n) \end{bmatrix}
\end{equation}

摩擦项的界限可以假设具有形式
\begin{equation}
||F_v\dot{q}|| \leq v||\dot{q}||, \quad ||F_d(\dot{q})|| \leq k
\end{equation}
其中$v$和$k$对特定机械臂是已知的，$||\cdot||$是合适的范数。

可以包含在$F(\dot{q})$中的另一个摩擦项是静摩擦，其分量具有形式
\begin{equation}
F_{si} = k_{si}\text{sgn}(\dot{q}_i), \quad |\dot{q}_i| < \epsilon
\end{equation}
其中$k_{si}$是关节$i$的静摩擦系数，$\epsilon$是一个小的正参数。我们通常忽略这个项。

\textbf{干扰项$\tau_d$的性质}

机械臂方程(3.3.1)有一个干扰项$\tau_d$，它可以表示建模不准确的动力学等。我们假设它是有界的，因此
\begin{equation}
||\tau_d|| \leq d
\end{equation}
其中$d$是可以为给定机械臂计算的标量常数，$||\cdot||$是任何合适的范数。

\subsection{参数线性}

机器人动力学方程享有的最后一个性质将在第五章中对我们非常有用。即，它在参数中是线性的，这个性质首先在[Craig 1988]的自适应控制中被利用。这很重要，因为一些或所有参数可能是未知的；因此动力学在未知项中是线性的。

这个性质可以表示为
\begin{equation}
M(q)\ddot{q} + V(q,\dot{q}) + F_v\dot{q} + F_d(\dot{q}) + G(q) = W(q,\dot{q},\ddot{q})\phi
\end{equation}
其中$\phi$是参数向量，$W(q,\dot{q},\ddot{q})$是取决于关节变量、关节速度和关节加速度的机器人函数矩阵。这个矩阵可以为任何给定机器人机械臂计算，因此是已知的。见例3.3.1。注意干扰$\tau_d$不包含在这个方程中。

\textbf{例3.3-1：二连杆平面肘形机械臂的结构和界限}

二连杆平面机械臂的动力学在例3.2.2中给出。我们现在应该计算本节中定义的结构矩阵，以及表3.3.1中需要的界限。摩擦界限是直接的，因此我们在这里不提及它们。动力学矩阵为
\begin{align}
M(q) &= \begin{bmatrix} (m_1+m_2)a_1^2 + m_2a_2^2 + 2m_2a_1a_2\cos\theta_2 & m_2a_2^2 + m_2a_1a_2\cos\theta_2 \\ m_2a_2^2 + m_2a_1a_2\cos\theta_2 & m_2a_2^2 \end{bmatrix} \\
V(q,\dot{q}) &= \begin{bmatrix} -m_2a_1a_2(2\dot{\theta}_1+\dot{\theta}_2)\dot{\theta}_2\sin\theta_2 \\ m_2a_1a_2\dot{\theta}_1^2\sin\theta_2 \end{bmatrix} \\
G(q) &= \begin{bmatrix} (m_1+m_2)ga_1\cos\theta_1 + m_2ga_2\cos(\theta_1+\theta_2) \\ m_2ga_2\cos(\theta_1+\theta_2) \end{bmatrix}
\end{align}

在表3.3.1中选择合适的范数并不总是简单的。在后续章节中要开发的控制算法中，我们证明了就某种范数而言的适当性能，这通常可以是任何期望的范数。为了实现控制器，必须选择特定的范数并评估界限。这个选择通常简单地取决于哪种范数使得可以评估表中的界限。例如，选择向量的2-范数需要评估$M(q)$的最大奇异值，这是一个非常困难的任务。

选择向量的$\infty$-范数意味着在每个采样时间确定$V(q(t),\dot{q}(t))$中具有最大幅度的元素。这需要决策逻辑，范数可能不是连续的。因此，让我们在这个例子中使用1-范数。相应的矩阵诱导范数则是最大绝对列和（第2章）。

\textbf{a. 惯量矩阵的界限}

$\mu_1$和$\mu_2$的评估相当于确定$M(q)$在所有$q$上的最小和最大特征值。这不是一件容易的事，需要求解一些二次方程，尽管使用Mathematica或Maple等软件可以相当容易地进行。因此，让我们找到$m_1$和$m_2$。

$M(q)$的诱导1-范数是最大绝对列和。在确定这个范数的界限时，重要的是要考虑关节角允许的运动范围。为了说明，假设$\theta_1$和$\theta_2$被限制在$\pm\pi/2$以内。那么1-范数总是以第1列表示为
\begin{equation}
||M(q)||_1 = |(m_1+m_2)a_1^2 + m_2a_2^2 + 2m_2a_1a_2\cos\theta_2| + |m_2a_2^2 + m_2a_1a_2\cos\theta_2|
\end{equation}

这被下式上界限制
\begin{equation}
m_2 = (m_1+m_2)a_1^2 + 2m_2a_2^2 + 3m_2a_1a_2
\end{equation}
和下界
\begin{equation}
m_1 = (m_1+m_2)a_1^2 + m_2a_2^2 - 2m_2a_1a_2
\end{equation}

由于机械臂是旋转的，且$\cos\theta_2$上下有界，$M_2$和$M_1$是常数。重要的是要注意，如果机械臂是旋转/平移（RP）的，因此关节变量为$(\theta_1, a_2)$，界限是$q$的函数。

\textbf{b. 科里奥利和重力项的界限}

使用
\begin{equation}
V(q,\dot{q}) = \begin{bmatrix} -m_2a_1a_2(2\dot{\theta}_1+\dot{\theta}_2)\dot{\theta}_2\sin\theta_2 \\ m_2a_1a_2\dot{\theta}_1^2\sin\theta_2 \end{bmatrix}
\end{equation}
找到科里奥利/向心向量的界限$v_b$，因此$v_b = m_2a_1a_2$。

同样，对于重力界限，
\begin{equation}
G(q) = \begin{bmatrix} (m_1+m_2)ga_1\cos\theta_1 + m_2ga_2\cos(\theta_1+\theta_2) \\ m_2ga_2\cos(\theta_1+\theta_2) \end{bmatrix}
\end{equation}

注意如果机械臂是RP，那么$v_b$和$g_b$是$q$的函数。

\textbf{c. 科里奥利/向心结构矩阵}

我们现在列出本节讨论的$V(q,\dot{q})$的各种结构矩阵。它们的计算留作练习（见习题）。

科里奥利/向心矩阵：
\begin{align}
V_{m1}(q,\dot{q}) &= \begin{bmatrix} -m_2a_1a_2\dot{\theta}_2\sin\theta_2 & -m_2a_1a_2(\dot{\theta}_1+\dot{\theta}_2)\sin\theta_2 \\ m_2a_1a_2\dot{\theta}_1\sin\theta_2 & 0 \end{bmatrix} \\
V_{m2}(q,\dot{q}) &= \begin{bmatrix} -m_2a_1a_2\dot{\theta}_2\sin\theta_2 & m_2a_1a_2\dot{\theta}_1\sin\theta_2 \\ -m_2a_1a_2(\dot{\theta}_1+\dot{\theta}_2)\sin\theta_2 & 0 \end{bmatrix}
\end{align}

斜对称矩阵$S(q,\dot{q})$：
\begin{equation}
S = \begin{bmatrix} 0 & -m_2a_1a_2(2\dot{\theta}_1+\dot{\theta}_2)\sin\theta_2 \\ m_2a_1a_2(2\dot{\theta}_1+\dot{\theta}_2)\sin\theta_2 & 0 \end{bmatrix}
\end{equation}

对称矩阵$V_1(q)$、$V_2(q)$：
\begin{align}
V_1(q) &= \begin{bmatrix} 0 & -m_2a_1a_2\sin\theta_2 \\ -m_2a_1a_2\sin\theta_2 & -m_2a_1a_2\sin\theta_2 \end{bmatrix} \\
V_2(q) &= \begin{bmatrix} m_2a_1a_2\sin\theta_2 & 0 \\ 0 & 0 \end{bmatrix}
\end{align}

位置/速度分解矩阵：
\begin{align}
V_p(q) &= \begin{bmatrix} 0 & -m_2a_1a_2(2\dot{\theta}_1+\dot{\theta}_2)\sin\theta_2 \\ m_2a_1a_2\dot{\theta}_1\sin\theta_2 & 0 \end{bmatrix} \\
V_v(\dot{q}) &= \begin{bmatrix} -m_2a_1a_2\dot{\theta}_2\sin\theta_2 & -m_2a_1a_2\dot{\theta}_2\sin\theta_2 \\ 0 & 0 \end{bmatrix}
\end{align}

\textbf{d. 机器人函数参数矩阵$W$}

机器人动力学在参数中是线性的。对于自适应控制的目的，应该选择表3.3.1中的参数向量$\phi$，使其包含未知参数。

包括摩擦在内的动力学可以写为
\begin{equation}
\tau = W\phi
\end{equation}
其中
\begin{equation}
\phi = \begin{bmatrix} m_1 \\ m_2 \\ v_1 \\ k_1 \\ v_2 \\ k_2 \end{bmatrix}
\end{equation}
第二质量$m_2$包括负载的质量。这个和摩擦系数通常是未知的。因此，选择
\begin{equation}
\phi = \begin{bmatrix} m_1a_1^2 + m_2a_1^2 \\ m_2a_2^2 \\ m_2a_1a_2 \\ m_1a_1 + m_2a_1 \\ m_2a_2 \\ v_1 \\ v_2 \\ k_1 \\ k_2 \end{bmatrix}
\end{equation}

那么已知机器人函数的矩阵$W(q,\dot{q},\ddot{q})$变为
\begin{equation}
W = \begin{bmatrix} w_{11} & w_{12} & w_{13} & w_{14} & w_{15} & \dot{\theta}_1 & \text{sgn}(\dot{\theta}_1) & 0 & 0 \\ 0 & w_{22} & w_{23} & 0 & w_{25} & 0 & 0 & \dot{\theta}_2 & \text{sgn}(\dot{\theta}_2) \end{bmatrix}
\end{equation}
其中
\begin{align}
w_{11} &= \ddot{\theta}_1 \\
w_{12} &= \ddot{\theta}_1 + \ddot{\theta}_2 \\
w_{13} &= (2\ddot{\theta}_1 + \ddot{\theta}_2)\cos\theta_2 - (2\dot{\theta}_1\dot{\theta}_2 + \dot{\theta}_2^2)\sin\theta_2 \\
w_{14} &= g\cos\theta_1 \\
w_{15} &= g\cos(\theta_1 + \theta_2) \\
w_{22} &= \ddot{\theta}_1 + \ddot{\theta}_2 \\
w_{23} &= \ddot{\theta}_1\cos\theta_2 + \dot{\theta}_1^2\sin\theta_2 \\
w_{25} &= g\cos(\theta_1 + \theta_2)
\end{align}

读者应该验证，有了这些定义，动力学可以表示为$\tau = W\phi$。矩阵$W$由测量的关节位置及其速度和加速度计算得出。

\subsection{无源性与能量守恒}

表3.3.1中给出的机械臂动力学的"牛顿"形式掩盖了一些我们应该在这里探讨的重要物理性质[Koditschek 1984]、[Ortega and Spong 1988]、[Johansson 1990]、[Slotine and Li 1987]、[Slotine 1988]。注意动力学可以用斜对称矩阵$S(q,\dot{q})$写为
\begin{equation}
M(q)\ddot{q} + V_m(q,\dot{q})\dot{q} + G(q) = \tau
\end{equation}
其中摩擦和$\tau_d$被忽略。现在，以$K$为动能，
\begin{equation}
\dot{K} = \frac{1}{2}\dot{q}^T\dot{M}\dot{q} + \dot{q}^T M\ddot{q}
\end{equation}
因此(3.3.63)产生
\begin{equation}
\dot{K} = \frac{1}{2}\dot{q}^T(\dot{M} - 2V_m)\dot{q} + \dot{q}^T(\tau - G)
\end{equation}
或
\begin{equation}
\dot{K} = \frac{1}{2}\dot{q}^T S \dot{q} + \dot{q}^T(\tau - G)
\end{equation}

这是能量守恒的表述，右边表示来自净外力的功率输入。$S(q,\dot{q})$的斜对称性只不过是一个陈述，即虚拟力$S(q,\dot{q})$不做功。外力所做的功由下式给出
\begin{equation}
K(t) - K(0) = \int_0^t \dot{q}^T(\tau - G) dt - \frac{1}{2}\int_0^t \dot{q}^T S \dot{q} dt = \int_0^t \dot{q}^T(\tau - G) dt
\end{equation}

此时回忆机械臂从$\tau(t)$到$\dot{q}$的无源性（第1.5节），它仅仅说明机械臂不能创造能量。从控制的角度来看，无源系统不能变得不稳定。一些流行的控制方案（例如，标准计算扭矩，第3.4节）的问题是它们破坏了无源性性质，如果系统参数不完全已知或存在干扰，可能导致不稳定。基于无源性的设计确保闭环系统是无源的（见第4.3节、上面引用的参考文献和[Anderson 1989]）。

这个分析不包括摩擦项。无论$f(\dot{q})$的形式如何，一个合理的假设是摩擦是耗散的，因此$f_i(x)$仅位于第一和第三象限。这等价于
\begin{equation}
\dot{q}^T F(\dot{q}) \geq 0
\end{equation}

在这一假设下，摩擦不会破坏机械臂的无源性。然后可以简单地修改(3.3.63)设计的控制器以包含摩擦[Slotine 1988]。摩擦的耗散性质允许人们将系统的带宽增加到经典极限之外。

\section{状态变量表示与反馈线性化}

机器人机械臂动力学方程在表3.3.1中为
\begin{equation}
M(q)\ddot{q} + V(q,\dot{q}) + F_v\dot{q} + F_d(\dot{q}) + G(q) + \tau_d = \tau
\end{equation}
其中$q(t) \in R^n$为关节变量向量，$\tau(t)$为控制输入。$M(q)$是惯量矩阵，$V(q,\dot{q})$是科里奥利/向心向量，$F_v(\dot{q})$是粘性摩擦，$F_d(\dot{q})$是动态摩擦，$G(q)$是重力，$\tau_d$是干扰。这些项满足表3.3.1中所示的性质。我们也可以将动力学写为
\begin{equation}
M(q)\ddot{q} + N(q,\dot{q}) = \tau
\end{equation}
其中非线性项表示为
\begin{equation}
N(q,\dot{q}) = V(q,\dot{q}) + F_v\dot{q} + F_d(\dot{q}) + G(q) + \tau_d
\end{equation}

在本节中，我们打算展示机械臂动力学方程的一些等价表述。

第2章讨论的非线性状态变量表示具有许多从控制角度来看有用的性质。函数$u(t)$是控制输入，$x(t)$是状态向量，描述能量在系统中如何存储。我们在这里展示如何将(3.4.1)放入这种形式。在第4章中，我们展示如何使用这种非线性状态变量形式使用计算机仿真机器人臂的行为。在整本书中，我们将反复使用状态空间表述进行控制设计，无论是非线性形式还是线性形式
\begin{equation}
\dot{x} = Ax + Bu
\end{equation}

在本节中，我们还介绍了非线性机器人方程反馈线性化的一般方法，它涉及以系统的方式重新定义变量，以产生关于我们感兴趣的动态变量的线性状态方程。这个变量可以是，例如，关节变量$q(t)$、笛卡尔位置，或相机参考系中的位置。

\subsection{哈密顿公式}

机械臂方程是使用拉格朗日力学推导的。在这里，让我们使用哈密顿力学[Marion 1965]推导机械臂动力学的状态变量表述[Arimoto and Miyazaki 1984]、[Gu and Loh 1985]。为了简单起见，让我们忽略摩擦项和干扰$\tau_d$；它们可以很容易地在我们的推导结束时添加。

在3.2节中，我们将机械臂拉格朗日量表示为
\begin{equation}
L(q,\dot{q}) = K(q,\dot{q}) - P(q) = \frac{1}{2}\dot{q}^T M(q)\dot{q} - P(q)
\end{equation}
其中$q(t) \in R^n$为关节变量，$K$为动能，$P$为势能，$M(q)$为机械臂惯量矩阵。定义广义动量为
\begin{equation}
p = \frac{\partial L}{\partial \dot{q}} = M(q)\dot{q}
\end{equation}

那么我们有
\begin{equation}
\dot{q} = M^{-1}(q)p
\end{equation}
以$p(t)$表示的动能为
\begin{equation}
K(q,p) = \frac{1}{2}p^T M^{-1}(q)p
\end{equation}

值得注意
\begin{equation}
\frac{\partial K}{\partial q} = -\frac{1}{2}(I_n \otimes p^T)\frac{\partial M^{-1}}{\partial q}p = \frac{1}{2}(I_n \otimes p^T)M^{-1}\frac{\partial M}{\partial q}M^{-1}p
\end{equation}

通过定义机械臂哈密顿量为
\begin{equation}
H(q,p) = K(q,p) + P(q)
\end{equation}
哈密顿运动方程为
\begin{align}
\dot{q} &= \frac{\partial H}{\partial p} = M^{-1}(q)p \\
\dot{p} &= -\frac{\partial H}{\partial q} + \tau = -\frac{\partial K}{\partial q} - \frac{\partial P}{\partial q} + \tau
\end{align}

注意
\begin{equation}
\frac{\partial P}{\partial q} = G(q)
\end{equation}

评估(3.4.13)得到
\begin{equation}
\dot{p} = \frac{1}{2}(I_n \otimes p^T)M^{-1}\frac{\partial M}{\partial q}M^{-1}p - G(q) + \tau
\end{equation}
可以表示为（见习题）
\begin{equation}
\dot{p} = -V_m(q,\dot{q})\dot{q} - G(q) + \tau
\end{equation}
其中$G(q)$是重力向量，$\otimes$是克罗内克积（见第3.3节）。

将状态向量定义为$x \in R^{2n}$
\begin{equation}
x = \begin{bmatrix} q \\ p \end{bmatrix}
\end{equation}
我们看到机械臂动力学可以表示为
\begin{equation}
\dot{x} = \begin{bmatrix} 0 & I_n \\ 0 & -V_m M^{-1} \end{bmatrix}x + \begin{bmatrix} 0 \\ I_n \end{bmatrix}u - \begin{bmatrix} 0 \\ G \end{bmatrix}
\end{equation}
其中控制输入定义为
\begin{equation}
u = \tau
\end{equation}

这是形式为(3.4.4)的非线性状态方程。重要的是要注意，这个动力学方程在控制输入$u$中是线性的，它激发广义动量$p(t)$的每个分量。

这个哈密顿状态空间表述在[Arimoto and Miyazaki 1984]中使用李雅普诺夫方法推导PID控制律，在[Gu and Loh 1985]中推导轨迹跟踪控制。

\subsection{位置/速度公式}

通过将位置/速度状态$x \in R^{2n}$定义为
\begin{equation}
x = \begin{bmatrix} q \\ \dot{q} \end{bmatrix}
\end{equation}
可以获得机械臂动力学的替代状态空间表述。

为了简单起见，忽略干扰$\tau_d$和摩擦$F_v + F_d(\dot{q})$，注意根据(3.4.2)，我们可以写出
\begin{equation}
\ddot{q} = M^{-1}(q)[\tau - N(q,\dot{q})]
\end{equation}

现在，我们可以直接写出位置/速度状态空间表示
\begin{equation}
\dot{x} = \begin{bmatrix} \dot{q} \\ -M^{-1}N \end{bmatrix} + \begin{bmatrix} 0 \\ M^{-1} \end{bmatrix}\tau = f(x) + B(x)\tau
\end{equation}
这是(3.4.4)的形式，$u(t) = \tau(t)$。

也可以写出一个形式为(3.4.5)的替代线性状态方程
\begin{equation}
\dot{x} = \begin{bmatrix} 0 & I \\ 0 & 0 \end{bmatrix}x + \begin{bmatrix} 0 \\ I \end{bmatrix}u
\end{equation}
其中控制输入定义为
\begin{equation}
u = M^{-1}(q)[\tau - N(q,\dot{q})]
\end{equation}

这两个位置/速度状态空间表述在后续章节中都将证明是有用的。

\subsection{反馈线性化}

现在让我们开发一种通用的方法，用于确定机械臂动力学(3.4.1)-(3.4.2)的线性状态空间表示。该技术涉及一种线性化变换，它消除机械臂的非线性。它是[Hunt et al. 1983, Gilbert and Ha 1984]中反馈线性化技术的简化版本。另见[Kreutz 1989]。

机器人动力学由(3.4.2)给出，$q \in R^n$。让我们定义一种通用类型的输出
\begin{equation}
y = h(q) - s(t)
\end{equation}
其中$h(q)$是关节变量$q \in R^n$的通用预定函数，$s(t)$是通用预定时间函数。那么控制问题将是选择关节扭矩和力输入$\tau(t)$以使输出$y(t)$变为零。

$h(q)$和$s(t)$的选择基于我们心中的控制目标。例如，如果$h(q) = -q$和$s(t) = q_d(t)$，我们希望机械臂遵循的期望关节空间轨迹，那么$y(t) = q_d(t) - q(t) = e(t)$，关节空间跟踪误差。在这种情况下强制$y(t)$为零将使关节变量$q(t)$跟踪它们的期望值$q_d(t)$，导致机械臂轨迹跟踪。

作为另一个例子，$h(q)$可以表示笛卡尔空间跟踪误差，$y_d$为位置误差，$e_0 \in R^3$为方向误差。将$y(t)$控制为零然后将导致直接在笛卡尔空间中的轨迹跟踪，毕竟，期望运动通常是在那里指定的。

最后，$-h(q)$可以表示到相机坐标系的非线性变换，$s(t)$是该坐标系中的期望轨迹。那么$y(t)$是相机坐标系跟踪误差。强制$y(t)$为零将产生相机空间中的跟踪运动。

\textbf{反馈线性化变换}

为了确定用于机器人控制器设计的线性状态变量模型，让我们简单地对输出$y(t)$求导两次以获得
\begin{equation}
\dot{y} = \frac{\partial h}{\partial q}\dot{q} - \dot{s} = J\dot{q} - \dot{s}
\end{equation}
\begin{equation}
\ddot{y} = J\ddot{q} + \dot{J}\dot{q} - \ddot{s}
\end{equation}
其中我们定义了雅可比矩阵
\begin{equation}
J(q) \equiv \frac{\partial h}{\partial q}
\end{equation}

如果$y \in R^p$，雅可比矩阵是$p \times n$形式的矩阵
\begin{equation}
J = \begin{bmatrix} \frac{\partial h_1}{\partial q_1} & \cdots & \frac{\partial h_1}{\partial q_n} \\ \vdots & \ddots & \vdots \\ \frac{\partial h_p}{\partial q_1} & \cdots & \frac{\partial h_p}{\partial q_n} \end{bmatrix}
\end{equation}

给定函数$h(q)$，计算与$h(q)$相关的雅可比矩阵$J(q)$是直接的。在特殊情况下，如果$\dot{q}$表示笛卡尔速度，$J(q)$是附录A中讨论的机械臂雅可比矩阵。那么，如果所有关节都是旋转的，$J$的单位是长度单位。

根据(3.4.2)，
\begin{equation}
\ddot{q} = M^{-1}(\tau - N)
\end{equation}
因此(3.4.26)产生
\begin{equation}
\ddot{y} = JM^{-1}(\tau - N) + \dot{J}\dot{q} - \ddot{s}
\end{equation}

定义控制输入函数
\begin{equation}
u = JM^{-1}(\tau - N) + \dot{J}\dot{q} - \ddot{s}
\end{equation}
和干扰函数
\begin{equation}
w = JM^{-1}\tau_d
\end{equation}

现在我们可以将状态$x(t) \in R^{2p}$定义为
\begin{equation}
x = \begin{bmatrix} y \\ \dot{y} \end{bmatrix}
\end{equation}
并将机器人动力学写为
\begin{equation}
\dot{x} = \begin{bmatrix} 0 & I \\ 0 & 0 \end{bmatrix}x + \begin{bmatrix} 0 \\ I \end{bmatrix}u + \begin{bmatrix} 0 \\ -w \end{bmatrix} = Ax + Bu + \begin{bmatrix} 0 \\ -w \end{bmatrix}
\end{equation}

这是形式为
\begin{equation}
\dot{x} = Ax + Bu + \begin{bmatrix} 0 \\ -w \end{bmatrix}
\end{equation}
的线性状态空间系统
由控制输入$u(t)$和干扰$v(t)$驱动。由于$A$和$B$的特殊形式，该系统被称为布鲁诺夫斯基标准形式（第2章）。读者应该确定可控性矩阵以验证它总是从$u(t)$可控。

方程(3.4.31)被称为机器人动力学方程的线性化变换。我们可以反转这个变换以获得
\begin{equation}
\tau = M J^+ (u - \dot{J}\dot{q} + \ddot{s}) + N + J^+ w
\end{equation}
其中$J^+$是雅可比矩阵$J(q)$的摩尔-彭若斯逆[Rao and Mitra 1971]。如果$J(q)$是方阵（即$p = n$）且非奇异，那么$J^+(q) = J^{-1}(q)$，我们可以写出
\begin{equation}
\tau = M J^{-1} (u - \dot{J}\dot{q} + \ddot{s}) + N + J^{-1} w
\end{equation}

正如我们将在第4章看到的，反馈线性化提供了一种强大的控制设计技术。事实上，如果我们选择$u(t)$使得(3.4.34)稳定（例如，一种可能性是PD反馈$u = -K_v \dot{y} - K_p y$），那么由(3.4.36)定义的控制输入扭矩$\tau(t)$使机器人臂以$y(t)$变为零的方式移动。

在特殊情况$y(t) = q(t)$，那么$J = I$，(3.4.34)简化为线性位置/速度形式(3.4.22)。

\section{笛卡尔动力学与其他动力学}

在3.2节中，我们以$q(t)$的时间行为推导了机器人动力学。根据表3.3.1，
\begin{equation}
M(q)\ddot{q} + N(q,\dot{q}) = \tau
\end{equation}
或
\begin{equation}
M(q)\ddot{q} + V(q,\dot{q}) + F(\dot{q}) + G(q) + \tau_d = \tau
\end{equation}
其中非线性项为
\begin{equation}
N(q,\dot{q}) = V(q,\dot{q}) + F(\dot{q}) + G(q) + \tau_d
\end{equation}
我们称之为以关节空间形式表述的机械臂动力学，或简称为关节空间动力学。

\subsection{笛卡尔机械臂动力学}

拥有除关节变量$q(t)$以外的变量的动力学发展描述通常是有用的。因此，定义
\begin{equation}
y = h(q)
\end{equation}
其中$h(q)$通常是来自运动学的非线性变换。虽然$y(t)$可以是任何感兴趣的变量，让我们在这里将其视为末端执行器的笛卡尔或任务空间位置（即在基坐标中末端执行器的位置和方向）。

从关节空间动力学到笛卡尔动力学的推导类似于3.4节中的反馈线性化。将(3.5.4)求导两次得到
\begin{align}
\dot{y} &= \frac{\partial h}{\partial q}\dot{q} = J\dot{q} \\
\ddot{y} &= J\ddot{q} + \dot{J}\dot{q}
\end{align}
其中雅可比矩阵为
\begin{equation}
J = \frac{\partial h}{\partial q}
\end{equation}

笛卡尔速度向量为$\dot{y} = [v^T \ \omega^T]^T$，其中$v \in R^3$是线速度，$\omega \in R^3$是角速度。让我们假设连杆数为$n = 6$，因此$J$是方阵。还假设我们远离工作空间奇异点，因此$|J| \neq 0$，根据(3.5.6)，我们可以写出
\begin{equation}
\ddot{q} = J^{-1}\ddot{y} - J^{-1}\dot{J}\dot{q}
\end{equation}
这是"逆加速度"变换。将其代入(3.5.2)得到
\begin{equation}
MJ^{-1}\ddot{y} - MJ^{-1}\dot{J}\dot{q} + N = \tau
\end{equation}

现在回忆力变换$\tau = J^T F$，其中$F$是笛卡尔力向量（见附录A），我们有
\begin{equation}
MJ^{-1}\ddot{y} - MJ^{-1}\dot{J}J^{-1}\dot{y} + N = J^T F
\end{equation}

这可以写为
\begin{equation}
M'\ddot{y} + N' = F
\end{equation}
其中我们定义了笛卡尔惯量矩阵、非线性项和干扰为
\begin{align}
M'(q) &= J^{-T}M(q)J^{-1} \\
N'(q,\dot{q}) &= J^{-T}[N(q,\dot{q}) - M(q)J^{-1}\dot{J}\dot{q}] \\
\tau_d' &= J^{-T}\tau_d
\end{align}

方程(3.5.9)-(3.5.10)给出了机器人机械臂的笛卡尔或工作空间动力学。

注意，$M'$、$N'$和$\tau_d'$取决于$q$和$\dot{q}$，因此严格来说，笛卡尔动力学并没有完全以$y$给出。然而，给定$y(t)$，我们可以使用逆运动学确定$q(t)$，因此$M'$、$N'$、$\tau_d'$可以作为$y$和$\dot{y}$的函数使用计算机子程序计算。

\subsection{笛卡尔动力学的结构与性质}

重要的是要认识到，只要$J$是非奇异的，表3.3.1中列出的关节空间动力学的所有性质都适用于笛卡尔动力学[Slotine and Li 1987]。特别要注意，$M'$是对称且正定的。对于旋转机械臂，雅可比矩阵具有长度单位且是有界的。在这种情况下，$M'$上下有界。

定义
\begin{equation}
V' = J^{-T}(V - MJ^{-1}\dot{J}J^{-1}\dot{y}) = J^{-T}(V - MJ^{-1}\dot{J}\dot{q})
\end{equation}
它遵循
\begin{equation}
V' = V_m'\dot{y}
\end{equation}
其中
\begin{equation}
V_m' = J^{-T}V_m J^{-1} - J^{-T}M J^{-1}\dot{J} J^{-1}
\end{equation}
其中$V_m$在第3.3节中定义。

很容易证明
\begin{equation}
S' \equiv \dot{M}' - 2V_m'
\end{equation}
是斜对称的。确实，使用恒等式
\begin{equation}
\frac{dJ^{-1}}{dt} = -J^{-1}\dot{J}J^{-1}
\end{equation}
我们可以看到
\begin{equation}
S' = J^{-T}SJ^{-1}
\end{equation}
这是斜对称的，因为$S$是斜对称的。

笛卡尔动力学中的摩擦项为
\begin{equation}
F' = J^{-T}F(\dot{q}) = J^{-T}F(J^{-1}\dot{y}) \equiv F'(\dot{y})
\end{equation}
它们满足表3.3.1中那样的界限。注意在笛卡尔坐标中，摩擦效应不是解耦的（例如，$J^{-T}F_vJ^{-1}$不是对角矩阵）。笛卡尔重力向量
\begin{equation}
G' = J^{-T}G(q)
\end{equation}
是有界的。

参数线性性质成立并表示为
\begin{equation}
M'\ddot{y} + V' + F' + G' = W'(y,\dot{y},\ddot{y})\phi
\end{equation}
其中已知笛卡尔机器人函数为
\begin{equation}
W' = J^{-T}W
\end{equation}
$\phi$是机械臂参数向量。

\textbf{例3.5-1：三连杆圆柱形机械臂的笛卡尔动力学}

让我们展示如何将例3.2.3中找到的关节空间动力学转换为笛卡尔动力学。从例A.3-1，机械臂雅可比矩阵为
\begin{equation}
J = \begin{bmatrix} -r\sin\theta & 0 & \cos\theta \\ r\cos\theta & 0 & \sin\theta \\ 0 & 1 & 0 \end{bmatrix}
\end{equation}
因此其逆为
\begin{equation}
J^{-1} = \begin{bmatrix} -\frac{\sin\theta}{r} & \frac{\cos\theta}{r} & 0 \\ 0 & 0 & 1 \\ \cos\theta & \sin\theta & 0 \end{bmatrix}
\end{equation}

从例3.2.3，机械臂惯量矩阵为
\begin{equation}
M = \begin{bmatrix} J + m_1r^2 + m_2r^2 & 0 & 0 \\ 0 & m_1 + m_2 & 0 \\ 0 & 0 & m_2 \end{bmatrix}
\end{equation}

应用(3.5.11)得到（验证！）
\begin{equation}
M' = \begin{bmatrix} \frac{J+m_1r^2+m_2r^2}{r^2} & 0 & 0 \\ 0 & m_1+m_2 & 0 \\ 0 & 0 & m_2 \end{bmatrix}
\end{equation}
其中$J' = \frac{J}{r^2} + m_1 + m_2$。

以类似的方式，可以计算$N'$。

\section{执行器动力学}

我们已经在关节空间坐标和笛卡尔坐标中讨论了刚性机器人机械臂的动力学。然而，机器人需要执行器来移动它；这些通常是电动或液压马达。因此，现在需要将执行器动力学添加到机械臂动力学中，以获得机械臂加执行器系统的完整动力学描述。关于执行器和传感器的良好参考文献由[de Silva 1989]提供。

\subsection{带执行器的机械臂动力学}

我们将考虑电动执行器的情况，假设电机是电枢控制的。液压执行器由类似的方程描述。在本小节中，我们假设电枢电感可以忽略。

$n$-连杆机器人机械臂的方程来自表3.3.1：
\begin{equation}
M(q)\ddot{q} + V(q,\dot{q}) + F(\dot{q}) + G(q) + \tau_d = \tau
\end{equation}
其中$q \in R^n$是机械臂关节变量。驱动连杆的电枢控制直流电机的动力学由$n$个解耦方程给出
\begin{equation}
J_M\ddot{q}_M + B_M\dot{q}_M + F_M + \tau = K_M v
\end{equation}
其中$\dot{q}_M = \text{vec}\{\dot{q}_{Mi}\}$，$\dot{q}_{Mi}$是第$i$个转子位置角，$\text{vec}\{\alpha_i\}$表示分量为$\alpha_i$的向量。控制输入是电机电压向量$v \in R^n$。

执行器系数矩阵都是常数，由下式给出
\begin{equation}
J_M = \text{diag}\{J_{Mi}\}, \quad B_M = \text{diag}\{B_{Mi}\}, \quad K_M = \text{diag}\{K_{Mi}R_{ai}^{-1}\}
\end{equation}
其中第$i$个电机具有惯量$J_{Mi}$、转子阻尼常数$B_{Mi}$、反电动势常数$K_{bi}$、扭矩常数$K_{Mi}$和电枢电阻$R_{ai}$。

从第$i$个电机到第$i$个机械臂连杆的耦合齿轮比为$r_i$，我们定义使得
\begin{equation}
q_i = r_i q_{Mi} \quad \text{或} \quad q = Rq_M
\end{equation}
如果第$i$个关节是旋转的，那么$r_i$是小于1的无量纲常数。如果$q_i$是平移的，那么$r_i$具有m/rad单位。

执行器摩擦向量由下式给出
\begin{equation}
F_M = \text{vec}\{F_{Mi}\}
\end{equation}
其中$F_{Mi}$是第$i$个转子的摩擦。

注意大写的"M"表示电机常数和变量，而$V_m$是用克里斯托费尔符号定义的机械臂科里奥利/向心向量。

使用(3.6.4)在(3.6.2)中消去$q_M$，然后代入(3.6.1)中的$\tau$，得到以关节变量表示的动力学
\begin{equation}
(M + R^2J_M)\ddot{q} + (V + R^2B_M)\dot{q} + (F + RF_M) + G + \tau_d = RK_M v
\end{equation}
或，通过适当定义符号，
\begin{equation}
M'(q)\ddot{q} + V'(q,\dot{q}) + F'(\dot{q}) + G(q) + \tau_d' = K'v
\end{equation}

\textbf{完整机械臂加执行器动力学的性质}

完整动力学(3.6.6)与机器人动力学(3.6.1)具有相同的形式。很容易验证完整机械臂加执行器动力学享有与表3.3.1中列出的机械臂动力学相同的性质（见习题）。特别地，$V'$是$\dot{M}'$和斜对称矩阵之间差的一半，所有有界性假设都成立，参数线性成立。因此，在我们设计控制器时，我们可以假设执行器已经包含在表3.3.1中的机械臂方程中。

\textbf{独立关节动力学}

在许多商业机器人机械臂中，齿轮比$r_i$非常小，在执行器/连杆耦合中提供很大的扭矩优势。这对机器人机械臂控制器的设计有重要的影响，大大简化了设计。

为了探索这一点，让我们按分量写出完整动力学
\begin{equation}
(J_{Mi}r_i^2 + m_{ii})\ddot{q}_i + (B_{Mi}r_i^2 + v_i)\dot{q}_i + k_i\text{sgn}(\dot{q}_i) + d_i = k_i'v_i
\end{equation}
其中$B \equiv \text{diag}\{B_i\}$，$d_i$是由下式给出的干扰
\begin{equation}
d_i = \sum_{j \neq i}m_{ij}\ddot{q}_j + \sum_{j,k}V_{jki}\dot{q}_j\dot{q}_k + g_i + f_i(\dot{q}_i)
\end{equation}
其中$m_{ij}$是$M'$的非对角元素，$V_{jki}$是第$i$个连杆摩擦的张量分量，$G_i$是第$i$个重力分量。

这个方程表明，如果$r_i$很小，机械臂动力学近似由$n$个具有常数系数的解耦二阶方程给出。关节耦合和重力的动力学效应仅作为干扰项乘以$r_i^2$出现。也就是说，机器人控制设计实际上只是控制执行器动力学的问题。

不幸的是，现代高性能任务使科里奥利和向心项变大，因此$d_i$不小。此外，现代高性能机械臂具有接近统一的齿轮比（例如，直驱机械臂），因此非线性必须在任何认真的控制设计中考虑。

\subsection{三阶机械臂加执行器动力学}

有时在控制设计中使用完整机器人机械臂的替代模型[Tarn et al. 1991]。它是一个三阶微分方程，当电机电枢电感不可忽略时应该使用。

当电枢电感$L_i$不可忽略时，我们必须使用电枢控制直流电机方程代替(3.6.2)
\begin{equation}
L\frac{dI}{dt} + RI + K_b'\dot{q}_M = v
\end{equation}
\begin{equation}
J_M\ddot{q}_M + B_M\dot{q}_M + F_M + \tau = K_T I
\end{equation}
其中$I \in R^n$是电枢电流向量，
\begin{equation}
L = \text{diag}\{L_i\}, \quad R = \text{diag}\{R_i\}, \quad K_T = \text{diag}\{K_{Ti}\}, \quad K_b' = \text{diag}\{K_{bi}r_i^{-1}\}
\end{equation}

重要的是要注意，$T$是电机电气时间常数矩阵。在前一小节中，假设这些时间常数与电机机械时间常数相比可以忽略不计。

为了确定机械臂加直流电机执行器的整体动力学，在(3.6.1)和(3.6.10)之间消去$\tau$以获得$I$的表达式。然后，求导以显式暴露$\dot{I}$。将这些表达式代入(3.6.9)（见习题）以获得形式为
\begin{equation}
L'M'(q)\dddot{q} + D(q,\dot{q},\ddot{q}) = v
\end{equation}
的动力学。

系数矩阵$D$由下式给出
\begin{equation}
D(q) = TM'(q)
\end{equation}
因此当$L_i$很小时可以忽略不计。

\subsection{具有关节柔性的动力学}

我们假设执行器和机器人连杆之间的耦合通过具有齿轮比$r_i$的刚性齿轮系提供。在实际实践中，耦合存在齿隙和齿轮系柔性或弹性。这里我们在机械臂动力学模型中包含关节的柔性，为了简单起见，假设$r_i = 1$。

这并不难做到。确实，假设耦合柔性被建模为刚性弹簧。那么在方程(3.6.1)、(3.6.2)中提到的扭矩只不过是
\begin{equation}
\tau = B_s(\dot{q}_M - \dot{q}) + K_s(q_M - q)
\end{equation}
其中$B_s = \text{diag}\{b_{si}\}$，$K_s = \text{diag}\{k_{si}\}$，$b_{si}$和$k_{si}$是第$i$个齿轮系的阻尼和弹簧常数。因此动力学方程变为
\begin{equation}
M(q)\ddot{q} + V(q,\dot{q}) + G(q) + B_s(\dot{q} - \dot{q}_M) + K_s(q - q_M) = 0
\end{equation}
\begin{equation}
J_M\ddot{q}_M + B_M\dot{q}_M + F_M - B_s(\dot{q}_M - \dot{q}) - K_s(q_M - q) = K_M v
\end{equation}

这些方程的结构与表3.3.1中描述的刚性关节机械臂非常不同。我们在第6章讨论具有关节柔性的机器人机械臂的控制（见[Spong 1987]）。下一个例子展示了在控制柔性关节机器人时可能出现的问题。

\textbf{例3.6-1：带柔性耦合轴的直流电机}

为了关注关节柔性的影响，让我们检查单个电枢控制的直流电机通过具有显著柔性的轴耦合到负载。电气和机械子系统如图3.6.1所示。

电机电气方程为
\begin{equation}
L\frac{di}{dt} + Ri + K_b\omega_m = u
\end{equation}
其中$i(t)$、$u(t)$分别是电枢电流和电压。反电动势为$e_b = K_b\omega_m$。

柔性轴施加的相互作用力由下式给出$f = b(\omega_m - \omega_L) + k(\theta_m - \theta_L)$，其中轴阻尼和弹簧常数分别表示为$b$和$k$。因此运动方程可以写为
\begin{equation}
J_m\dot{\omega}_m + b_m\omega_m + b(\omega_m - \omega_L) + k(\theta_m - \theta_L) = K_T i
\end{equation}
\begin{equation}
J_L\dot{\omega}_L - b(\omega_m - \omega_L) - k(\theta_m - \theta_L) = 0
\end{equation}
其中下标$m$和$L$分别指电机参数和负载参数。假设负载惯量$J_L$是常数。其余符号的定义可以从前面的文本中推断。

为了将这些方程放入状态空间形式，将状态定义为
\begin{equation}
x = \begin{bmatrix} \theta_L \\ \theta_m \\ \omega_L \\ \omega_m \\ i \end{bmatrix}
\end{equation}
其中$\omega_m$、$\omega_L$是电机和负载的角速度。那么
\begin{equation}
\dot{x} = \begin{bmatrix} 0 & 0 & 1 & 0 & 0 \\ 0 & 0 & 0 & 1 & 0 \\ -\frac{k}{J_L} & \frac{k}{J_L} & -\frac{b}{J_L} & \frac{b}{J_L} & 0 \\ \frac{k}{J_m} & -\frac{k}{J_m} & \frac{b}{J_m} & -\frac{(b+b_m)}{J_m} & \frac{K_T}{J_m} \\ 0 & 0 & 0 & -\frac{K_b}{L} & -\frac{R}{L} \end{bmatrix}x + \begin{bmatrix} 0 \\ 0 \\ 0 \\ 0 \\ \frac{1}{L} \end{bmatrix}u
\end{equation}

\textbf{a. 刚性耦合轴}

如果耦合轴没有柔性，$\omega_m = \omega_L = \omega$，状态方程简化为（见习题）
\begin{equation}
\dot{x} = \begin{bmatrix} 0 & 1 & 0 \\ 0 & -\frac{(b_m+b_L)}{J} & \frac{K_T}{J} \\ 0 & -\frac{K_b}{L} & -\frac{R}{L} \end{bmatrix}x + \begin{bmatrix} 0 \\ 0 \\ \frac{1}{L} \end{bmatrix}u
\end{equation}
其中$x = [i \ \omega]^T$，$J = J_m + J_L$。将输出定义为电机速度得到
\begin{equation}
y = \begin{bmatrix} 0 & 0 & 1 \end{bmatrix}x
\end{equation}

传递函数计算为
\begin{equation}
H(s) = \frac{K_T}{(Ls+R)[(J_m+J_L)s+(b_m+b_L)] + K_T K_b}
\end{equation}

使用参数值$J_m = J_L = 0.1$ kg-m$^2$，$L = 0.5$ H，$b_m = 0.2$ N-m/rad/s，和$R = 5$ $\Omega$得到
\begin{equation}
H(s) = \frac{20}{(s+2.3)(s+8.7)}
\end{equation}
因此有两个实极点在$s = -2.3$，$s = -8.7$。

使用附录B中的程序TRESP进行仿真（见第3.3节）产生图3.6.2中所示的$\omega$阶跃响应。

\textbf{b. 非常柔性的耦合轴}

对应于非常柔性轴的耦合轴参数$k = 2$ N-m/rad和$b = 0.2$ N-m/rad/s。使用这些值，可以使用PC-MATLAB等软件获得两个传递函数
\begin{equation}
H_m(s) = \frac{20s(s^2+3s+44)}{(s+10)[(s+3.4)^2+5.6^2]}
\end{equation}
\begin{equation}
H_L(s) = \frac{40(s+10)}{(s+10)[(s+3.4)^2+5.6^2]}
\end{equation}

轴柔性模式有极点$s = -3.4 \pm j5.6$，因此阻尼比为$\zeta = 0.52$，自然频率为$\omega = 6.55$ rad/s。注意系统是边际稳定的，有一个极点在$s = 0$。由于极零点抵消，它是BIBO稳定的。

程序TRESP产生图3.6.3中所示的阶跃响应。有几点值得注意。最初，电机速度$\omega_m$比图3.6.2中上升得更快，因为轴柔性意味着只有转子转动惯量$J_m$最初影响速度。然后，当负载$J_L$通过轴耦合回电机时，$\omega_m$的增加速率减慢。还要注意，负载速度$\omega_L$由于轴中的柔性表现出大约0.1 s的延迟。

极其有趣的是要注意，轴柔性具有加速最慢电机实极点的效果[比较(8)和(9)]，因此$\omega_L$比刚性轴情况下更快地接近其稳态值。这是由于柔性轴的"甩动"作用。

轴动力学使得在没有某种专门设计的控制器的情况下，控制$\theta_L$（在机器人机械臂中对应于关节角$q_i$）非常困难。

\section{总结}

在本章中，我们为机器人控制系统的研究奠定了基础。在3.2节中使用拉格朗日力学，我们推导了一些将在整本书中用于演示设计的机械臂动力学。我们为任何串联连杆机械臂提供了通用机器人机械臂动力学的表达式。

在3.3节中，我们研究了机器人动力学的性质，如有界性、参数线性和斜对称性，这些在控制设计中是必需的。表3.3.1给出了我们发现的总结。我们使用了一种克罗内克积方法，该方法对机器人方程中各项之间的关系提供了很好的洞察。

现代控制系统设计中的一个重要形式是状态变量形式。在3.4节中，我们推导了机械臂动力学的几种状态空间形式，为后续章节中提供的几种设计技术奠定了基础。状态形式在机器人控制器的计算机仿真中也是有用的，如我们在第4节中看到的那样。

笛卡尔形式的动力学在3.5节中给出。驱动机器人机械臂连杆的执行器的动力学在3.6节中分析并包含。

\section{参考文献}

\begin{enumerate}
\item [Anderson 1989] Anderson, R.J., "Passive computed torque algorithms for robots," Proc. IEEE Conf. Decision Control, pp. 1638--1644, Dec. 1989.

\item [Arimoto and Miyazaki 1984] Arimoto, S., and F.Miyazaki, "Stability and robustness of PID feedback control for robot manipulators of sensory capability," Proc. First Int. Symp., pp. 783--799, MIT, Cambridge, MA, 1984.

\item [Asada and Slotine 1986] Asada, H., and J.-J.E.Slotine, \textit{Robot Analysis and Control}, New York: Wiley, 1986.

\item [Borisenko and Tarapov 1968] Borisenko, A.I., and I.E.Tarapov, \textit{Vector and Tensor Analysis with Applications}. Englewood Cliffs, NJ: Prentice Hall, 1968.

\item [Brewer 1978] Brewer, J.W., "Kronecker products and matrix calculus in system theory," \textit{IEEE Trans. Circuits Syst.}, vol. CAS-25, no. 9, pp. 772--781, Sept. 1978.

\item [Craig 1988] Craig, J.J., \textit{Adaptive Control of Mechanical Manipulators}. Reading, MA: Addison-Wesley, 1988.

\item [de Silva 1989] de Silva, C.W., \textit{Control Sensors and Actuators}. Englewood Cliffs, NJ: Prentice Hall, 1989.

\item [Gilbert and Ha 1984] Gilbert, E.G., and I.J.Ha, "An approach to nonlinear feedback control with applications to robotics," \textit{IEEE Trans. Syst. Man Cybern.}, vol. SMC-14, no. 6, pp. 879--884, Nov./Dec. 1984.

\item [Gu and Loh 1985] Gu, Y.-L., and N.K.Loh, "Dynamic model for industrial robots based on a compact Lagrangian formulation," Proc. IEEE Conf. Decision Control, pp. 1497--1501, 1985.

\item [Gu and Loh 1988] Gu, Y-L., and N.K.Loh, "Dynamic modeling and control by utilizing an imaginary robot model," \textit{IEEE J. Robot. Autom.}, vol. 4, no. 5, pp. 532--534, Oct. 1988.

\item [Hunt et al. 1983] Hunt, L.R., R.Su, and G.Meyer, "Global transformations of nonlinear systems," \textit{IEEE Trans. Autom. Control}, vol. AC-28, no. 1, pp. 24--31, Jan. 1983.

\item [Johansson 1990] Johansson, R., "Quadratic optimization of motion coordination and control," \textit{IEEE Trans. Autom. Control}, vol. 35, no. 11, pp. 1197--1208, Nov. 1990.

\item [Koditschek 1984] Koditschek, D., "Natural motion for robot arms," Proc. IEEE Conf. Decision Control, pp. 733--735, Dec. 1984.

\item [Kreutz 1989] Kreutz, K., "On manipulator control by exact linearization," \textit{IEEE Trans. Autom. Control}, vol. 34, no. 7, pp. 763--767, July 1989.

\item [Lee et al. 1983] Lee, C.S.G., R.C.Gonzalez, and K.S.Fu, \textit{Tutorial on Robotics}. New York: IEEE Press, 1983.

\item [Marion 1965] Marion, J.B., \textit{Classical Dynamics}. New York: Academic Press, 1965.

\item [Ortega and Spong 1988] Ortega, R., and Spong, M.W., "Adaptive motion control of rigid robots: a tutorial," Proc. IEEE Conf. Decision Control, pp. 1575--1584, Dec. 1988.

\item [Paul 1981] Paul, R.P., \textit{Robot Manipulators}. Cambridge, MA: MIT Press, 1981.

\item [Rao and Mitra 1971] Rao, C.R., and S.K.Mitra, \textit{Generalized Inverse of Matrices and Its Applications}. New York: Wiley, 1971.

\item [Schilling 1990] Schilling, R.J., \textit{Fundamentals of Robotics}. Englewood Cliffs, NJ: Prentice Hall, 1990.

\item [Slotine 1988] Slotine, J.-J.E., "Putting physics in control: the example of robotics," \textit{IEEE Control Syst. Mag.}, pp. 12--17, Dec. 1988.

\item [Slotine and Li 1987] Slotine, J.-J.E., and W.Li, "Adaptive strategies in constrained manipulation," Proc. IEEE Conf. Robot. Autom., pp. 595--601, 1987.

\item [Spong 1987] Spong, M.W., "Modeling and control of elastic joint robots," \textit{J. Dyn. Syst. Meas. Control}, vol. 109, pp. 310--319, Dec. 1987.

\item [Spong and Vidyasagar 1989] Spong, M.W., and M.Vidyasagar, \textit{Robot Dynamics and Control}. New York: Wiley, 1989.

\item [Tarn et al. 1991] Tarn, T.-J., A.K.Bejczy, X.Yun, and Z.Li, "Effect of motor dynamics on nonlinear feedback robot arm control," \textit{IEEE Trans. Robot. Autom.}, vol. 7, no. 1, pp. 114--122, Feb. 1991.
\end{enumerate}

\section{习题}

\textbf{第3.2节}

\begin{enumerate}
\item[3.2-1] \textbf{动力学}。求例A.2-4中球腕的动力学。

\item[3.2-2] \textbf{从推导方程的动力学}。在例3.2.2中，我们从基本原理找到了二连杆平面肘形机械臂的动力学。在本题中，从该例子中动能和势能的表达式开始，并：
\begin{enumerate}
\item[(a)] 将$K$写成形式(3.2.29)以确定$M(q)$。
\item[(b)] 使用(3.2.42)和(3.2.43)确定$V(q,\dot{q})$和$G(q)$。
\end{enumerate}

\item[3.2-3] \textbf{从推导方程的动力学}。对例3.2-3中的三连杆机械臂重复问题3.2-2。
\end{enumerate}

\textbf{第3.3节}

\begin{enumerate}
\item[3.3-1] 通过找到$V_{p1}(q)$和$V_{v1}(q)$证明(3.3.22)。

\item[3.3-2] 通过找到矩阵$V_i(q)$证明(3.3.23)。

\item[3.3-3] 证明(3.3.27)。

\item[3.3-4] \textbf{科里奥利项}。在(3.3.40)中找到$V_{cor}(q)$和$V_{cen}(q)$。

\item[3.3-5] \textbf{科里奥利项}。证明动力学方程中的科里奥利/向心项可以表示为$V(q,\dot{q}) = \text{vec}\{V(q,\dot{q})\}$，其中
\[ V = \sum_{i=1}^{n} \sum_{j=1}^{n} \sum_{k=1}^{n} \text{tr}\left(\frac{\partial T_i}{\partial q_j}M_i\frac{\partial T_i^T}{\partial q_k}\right)\dot{q}_j\dot{q}_k \]
其中
\[ \frac{\partial T_i}{\partial q_j} = T_i Q_j \]
和附录A中定义的$T_i$。将其与第3.2节中定义的$V_{m1}$、$V_{m2}$、$V_m$进行比较。

\item[3.3-6] \textbf{界限与结构}。详细推导例3.3.1中的结果。

\item[3.3-7] \textbf{界限与结构}。为例3.2.1中的二连杆极坐标机械臂推导界限和结构矩阵。使用：
\begin{enumerate}
\item[(a)] 1-范数。
\item[(b)] 2-范数。
\item[(c)] $\infty$-范数。
\end{enumerate}

\item[3.3-8] \textbf{界限与结构}。对习题3.2.3中的三连杆圆柱形机械臂重复问题3.3-7。

\item[3.3-9] \textbf{使用2-范数的界限}。使用2-范数推导例3.3.1中二连杆平面肘形机械臂的界限。
\end{enumerate}

\textbf{第3.4节}

\begin{enumerate}
\item[3.4-1] 证明(3.4.15)。

\item[3.4-2] \textbf{哈密顿状态公式}。证明(3.4.15)等价于
\[ \dot{p} = -V_m^T(q,\dot{q})\dot{q} - G(q) + \tau \]
其中斜对称矩阵在第3.3节中定义。

\item[3.4-3] \textbf{哈密顿状态公式}。使用(3.4.17)为例3.2.1中的二连杆极坐标机械臂推导哈密顿状态变量公式。

\item[3.4-4] \textbf{哈密顿状态公式}]。为例3.2.2中的二连杆平面肘形机械臂重复问题3.4-3。
\end{enumerate}

\textbf{第3.5节}

\begin{enumerate}
\item[3.5-1] \textbf{笛卡尔动力学}。完成例3.5.1，计算笛卡尔坐标中的非线性项。

\item[3.5-2] \textbf{笛卡尔动力学}。求例3.2.1中二连杆极坐标机械臂的笛卡尔动力学。

\item[3.5-3] \textbf{笛卡尔动力学}。求例3.2.2中二连杆平面肘形机械臂的笛卡尔动力学。
\end{enumerate}

\textbf{第3.6节}

\begin{enumerate}
\item[3.6-1] \textbf{执行器动力学}。验证机械臂加执行器动力学(3.6.6)具有表3.3.1中列出的性质。

\item[3.6-2] \textbf{执行器动力学}。推导三阶动力学(3.6.12)，提供$\tau$、$N$的显式表达式。验证当$L_i$可忽略时它们简化为(3.6.5)。

\item[3.6-3] \textbf{柔性耦合轴}。验证例3.6.1中刚性轴情况的状态方程。
\end{enumerate}

\end{document}
